% Options for packages loaded elsewhere
\PassOptionsToPackage{unicode}{hyperref}
\PassOptionsToPackage{hyphens}{url}
\documentclass[
]{article}
\usepackage{xcolor}
\usepackage[margin=1in]{geometry}
\usepackage{amsmath,amssymb}
\setcounter{secnumdepth}{-\maxdimen} % remove section numbering
\usepackage{iftex}
\ifPDFTeX
  \usepackage[T1]{fontenc}
  \usepackage[utf8]{inputenc}
  \usepackage{textcomp} % provide euro and other symbols
\else % if luatex or xetex
  \usepackage{unicode-math} % this also loads fontspec
  \defaultfontfeatures{Scale=MatchLowercase}
  \defaultfontfeatures[\rmfamily]{Ligatures=TeX,Scale=1}
\fi
\usepackage{lmodern}
\ifPDFTeX\else
  % xetex/luatex font selection
\fi
% Use upquote if available, for straight quotes in verbatim environments
\IfFileExists{upquote.sty}{\usepackage{upquote}}{}
\IfFileExists{microtype.sty}{% use microtype if available
  \usepackage[]{microtype}
  \UseMicrotypeSet[protrusion]{basicmath} % disable protrusion for tt fonts
}{}
\makeatletter
\@ifundefined{KOMAClassName}{% if non-KOMA class
  \IfFileExists{parskip.sty}{%
    \usepackage{parskip}
  }{% else
    \setlength{\parindent}{0pt}
    \setlength{\parskip}{6pt plus 2pt minus 1pt}}
}{% if KOMA class
  \KOMAoptions{parskip=half}}
\makeatother
\usepackage{color}
\usepackage{fancyvrb}
\newcommand{\VerbBar}{|}
\newcommand{\VERB}{\Verb[commandchars=\\\{\}]}
\DefineVerbatimEnvironment{Highlighting}{Verbatim}{commandchars=\\\{\}}
% Add ',fontsize=\small' for more characters per line
\usepackage{framed}
\definecolor{shadecolor}{RGB}{248,248,248}
\newenvironment{Shaded}{\begin{snugshade}}{\end{snugshade}}
\newcommand{\AlertTok}[1]{\textcolor[rgb]{0.94,0.16,0.16}{#1}}
\newcommand{\AnnotationTok}[1]{\textcolor[rgb]{0.56,0.35,0.01}{\textbf{\textit{#1}}}}
\newcommand{\AttributeTok}[1]{\textcolor[rgb]{0.13,0.29,0.53}{#1}}
\newcommand{\BaseNTok}[1]{\textcolor[rgb]{0.00,0.00,0.81}{#1}}
\newcommand{\BuiltInTok}[1]{#1}
\newcommand{\CharTok}[1]{\textcolor[rgb]{0.31,0.60,0.02}{#1}}
\newcommand{\CommentTok}[1]{\textcolor[rgb]{0.56,0.35,0.01}{\textit{#1}}}
\newcommand{\CommentVarTok}[1]{\textcolor[rgb]{0.56,0.35,0.01}{\textbf{\textit{#1}}}}
\newcommand{\ConstantTok}[1]{\textcolor[rgb]{0.56,0.35,0.01}{#1}}
\newcommand{\ControlFlowTok}[1]{\textcolor[rgb]{0.13,0.29,0.53}{\textbf{#1}}}
\newcommand{\DataTypeTok}[1]{\textcolor[rgb]{0.13,0.29,0.53}{#1}}
\newcommand{\DecValTok}[1]{\textcolor[rgb]{0.00,0.00,0.81}{#1}}
\newcommand{\DocumentationTok}[1]{\textcolor[rgb]{0.56,0.35,0.01}{\textbf{\textit{#1}}}}
\newcommand{\ErrorTok}[1]{\textcolor[rgb]{0.64,0.00,0.00}{\textbf{#1}}}
\newcommand{\ExtensionTok}[1]{#1}
\newcommand{\FloatTok}[1]{\textcolor[rgb]{0.00,0.00,0.81}{#1}}
\newcommand{\FunctionTok}[1]{\textcolor[rgb]{0.13,0.29,0.53}{\textbf{#1}}}
\newcommand{\ImportTok}[1]{#1}
\newcommand{\InformationTok}[1]{\textcolor[rgb]{0.56,0.35,0.01}{\textbf{\textit{#1}}}}
\newcommand{\KeywordTok}[1]{\textcolor[rgb]{0.13,0.29,0.53}{\textbf{#1}}}
\newcommand{\NormalTok}[1]{#1}
\newcommand{\OperatorTok}[1]{\textcolor[rgb]{0.81,0.36,0.00}{\textbf{#1}}}
\newcommand{\OtherTok}[1]{\textcolor[rgb]{0.56,0.35,0.01}{#1}}
\newcommand{\PreprocessorTok}[1]{\textcolor[rgb]{0.56,0.35,0.01}{\textit{#1}}}
\newcommand{\RegionMarkerTok}[1]{#1}
\newcommand{\SpecialCharTok}[1]{\textcolor[rgb]{0.81,0.36,0.00}{\textbf{#1}}}
\newcommand{\SpecialStringTok}[1]{\textcolor[rgb]{0.31,0.60,0.02}{#1}}
\newcommand{\StringTok}[1]{\textcolor[rgb]{0.31,0.60,0.02}{#1}}
\newcommand{\VariableTok}[1]{\textcolor[rgb]{0.00,0.00,0.00}{#1}}
\newcommand{\VerbatimStringTok}[1]{\textcolor[rgb]{0.31,0.60,0.02}{#1}}
\newcommand{\WarningTok}[1]{\textcolor[rgb]{0.56,0.35,0.01}{\textbf{\textit{#1}}}}
\usepackage{longtable,booktabs,array}
\newcounter{none} % for unnumbered tables
\usepackage{calc} % for calculating minipage widths
% Correct order of tables after \paragraph or \subparagraph
\usepackage{etoolbox}
\makeatletter
\patchcmd\longtable{\par}{\if@noskipsec\mbox{}\fi\par}{}{}
\makeatother
% Allow footnotes in longtable head/foot
\IfFileExists{footnotehyper.sty}{\usepackage{footnotehyper}}{\usepackage{footnote}}
\makesavenoteenv{longtable}
\usepackage{graphicx}
\makeatletter
\newsavebox\pandoc@box
\newcommand*\pandocbounded[1]{% scales image to fit in text height/width
  \sbox\pandoc@box{#1}%
  \Gscale@div\@tempa{\textheight}{\dimexpr\ht\pandoc@box+\dp\pandoc@box\relax}%
  \Gscale@div\@tempb{\linewidth}{\wd\pandoc@box}%
  \ifdim\@tempb\p@<\@tempa\p@\let\@tempa\@tempb\fi% select the smaller of both
  \ifdim\@tempa\p@<\p@\scalebox{\@tempa}{\usebox\pandoc@box}%
  \else\usebox{\pandoc@box}%
  \fi%
}
% Set default figure placement to htbp
\def\fps@figure{htbp}
\makeatother
\setlength{\emergencystretch}{3em} % prevent overfull lines
\providecommand{\tightlist}{%
  \setlength{\itemsep}{0pt}\setlength{\parskip}{0pt}}
\usepackage{bookmark}
\IfFileExists{xurl.sty}{\usepackage{xurl}}{} % add URL line breaks if available
\urlstyle{same}
\hypersetup{
  pdftitle={Cheatsheets},
  hidelinks,
  pdfcreator={LaTeX via pandoc}}

\title{Cheatsheets}
\author{}
\date{\vspace{-2.5em}}

\begin{document}
\maketitle

{
\setcounter{tocdepth}{3}
\tableofcontents
}
\section{Cheatsheets}\label{cheatsheets}

One-page topical summaries, aggregated here as a printable reference.
Each is a self-contained reminder of the syntax, the assumptions, and
the diagnostics that go with the methods it covers.

\subsection{Foundations}\label{foundations}

\subsubsection{\texorpdfstring{Toolchain, tidy data, and
\texttt{ggplot2}\#\# The research
workflow}{Toolchain, tidy data, and ggplot2\#\# The research workflow}}\label{toolchain-tidy-data-and-ggplot2-the-research-workflow}

\texttt{Question\ →\ Measurements\ →\ Design\ →\ Acquisition\ →\ Description\ →\ Analysis\ →\ Interpretation\ →\ Validation\ →\ Knowledge}

Biological vs statistical significance: ``\emph{p} \textless{} 0.05'' ≠
``effect matters''.

\subsection{R / RStudio / Quarto / renv}\label{r-rstudio-quarto-renv}

{\def\LTcaptype{none} % do not increment counter
\begin{longtable}[]{@{}
  >{\raggedright\arraybackslash}p{(\linewidth - 2\tabcolsep) * \real{0.5000}}
  >{\raggedright\arraybackslash}p{(\linewidth - 2\tabcolsep) * \real{0.5000}}@{}}
\toprule\noalign{}
\begin{minipage}[b]{\linewidth}\raggedright
Task
\end{minipage} & \begin{minipage}[b]{\linewidth}\raggedright
Command
\end{minipage} \\
\midrule\noalign{}
\endhead
\bottomrule\noalign{}
\endlastfoot
Install packages & \texttt{install.packages("pkg")} \\
Pin project & \texttt{renv::init()} / \texttt{renv::snapshot()} /
\texttt{renv::restore()} \\
Render one file & \texttt{quarto\ render\ path/to/file.qmd} \\
Preview with live reload & \texttt{quarto\ preview} \\
Render only slides &
\texttt{quarto\ render\ file.qmd\ -\/-to\ revealjs} \\
\end{longtable}
}

\subsection{Data types and scales}\label{data-types-and-scales}

{\def\LTcaptype{none} % do not increment counter
\begin{longtable}[]{@{}lll@{}}
\toprule\noalign{}
Scale & Example & Statistic \\
\midrule\noalign{}
\endhead
\bottomrule\noalign{}
\endlastfoot
Nominal & sex, site & counts, proportions \\
Ordinal & pain 0--10 & median, IQR, rank tests \\
Interval & °C & mean, SD \\
Ratio & kg, time & mean, SD, ratios \\
\end{longtable}
}

Accuracy ≠ precision. Bias shifts the target; variance spreads the
shots.

\subsection{Tidy data}\label{tidy-data}

\begin{enumerate}
\def\labelenumi{\arabic{enumi}.}
\tightlist
\item
  One variable per column.
\item
  One observation per row.
\item
  One observational unit per table.
\end{enumerate}

\subsection{\texorpdfstring{\texttt{dplyr}
verbs}{dplyr verbs}}\label{dplyr-verbs}

\begin{Shaded}
\begin{Highlighting}[]
\NormalTok{df }\SpecialCharTok{|\textgreater{}}
  \FunctionTok{filter}\NormalTok{(age }\SpecialCharTok{\textgreater{}=} \DecValTok{18}\NormalTok{) }\SpecialCharTok{|\textgreater{}}
  \FunctionTok{select}\NormalTok{(id, age, sbp) }\SpecialCharTok{|\textgreater{}}
  \FunctionTok{mutate}\NormalTok{(}\AttributeTok{sbp\_z =}\NormalTok{ (sbp }\SpecialCharTok{{-}} \FunctionTok{mean}\NormalTok{(sbp)) }\SpecialCharTok{/} \FunctionTok{sd}\NormalTok{(sbp)) }\SpecialCharTok{|\textgreater{}}
  \FunctionTok{group\_by}\NormalTok{(arm) }\SpecialCharTok{|\textgreater{}}
  \FunctionTok{summarise}\NormalTok{(}\AttributeTok{mean\_sbp =} \FunctionTok{mean}\NormalTok{(sbp, }\AttributeTok{na.rm =} \ConstantTok{TRUE}\NormalTok{), }\AttributeTok{n =} \FunctionTok{n}\NormalTok{())}
\end{Highlighting}
\end{Shaded}

Joins: \texttt{left\_join}, \texttt{inner\_join}, \texttt{anti\_join} by
a key.

Missingness: \texttt{is.na()}, \texttt{drop\_na()},
\texttt{tidyr::complete()}.

\subsection{\texorpdfstring{\texttt{ggplot2}
grammar}{ggplot2 grammar}}\label{ggplot2-grammar}

\begin{Shaded}
\begin{Highlighting}[]
\FunctionTok{ggplot}\NormalTok{(data, }\FunctionTok{aes}\NormalTok{(x, y, }\AttributeTok{colour =}\NormalTok{ group)) }\SpecialCharTok{+}
  \FunctionTok{geom\_point}\NormalTok{() }\SpecialCharTok{+}
  \FunctionTok{facet\_wrap}\NormalTok{(}\SpecialCharTok{\textasciitilde{}}\NormalTok{ factor) }\SpecialCharTok{+}
  \FunctionTok{labs}\NormalTok{(}\AttributeTok{x =} \StringTok{"x label"}\NormalTok{, }\AttributeTok{y =} \StringTok{"y label"}\NormalTok{) }\SpecialCharTok{+}
  \FunctionTok{theme\_minimal}\NormalTok{()}
\end{Highlighting}
\end{Shaded}

Layer order:
\texttt{data\ →\ aes\ →\ geom\ →\ facet\ →\ scales\ →\ labels\ →\ theme}.

\subsection{Plot families}\label{plot-families}

\begin{itemize}
\tightlist
\item
  One variable continuous → \textbf{histogram}, \textbf{density}.
\item
  Two continuous → \textbf{scatter}, sometimes binned / smoothed.
\item
  One categorical + one continuous → \textbf{boxplot}, \textbf{violin},
  \textbf{strip}.
\item
  Two categorical → \textbf{bar with fill}, \textbf{mosaic}.
\end{itemize}

\subsection{Decision rule}\label{decision-rule}

\begin{itemize}
\tightlist
\item
  Cannot read the data into a tibble → fix data entry first.
\item
  Tibble reads but has missing or mixed types → clean before modelling.
\item
  Plot is ugly or cluttered → trust the plot, not the summary statistic.
\end{itemize}

\subsection{Common pitfalls}\label{common-pitfalls}

\begin{itemize}
\tightlist
\item
  Running statistics on untyped / partially-missing data.
\item
  Adjusting colour scales instead of simplifying the chart.
\item
  Reporting mean + SD for a skewed variable.
\item
  Forgetting \texttt{set.seed()} in any simulation chunk.
\end{itemize}

\subsection{Further reading}\label{further-reading}

\begin{itemize}
\tightlist
\item
  Wickham, \emph{R for Data Science} (2e).
\item
  Posit ggplot2 and dplyr cheatsheets.
\end{itemize}

\subsubsection{Descriptive statistics, Bayes, and distributions\#\#
Descriptive
statistics}\label{descriptive-statistics-bayes-and-distributions-descriptive-statistics}

{\def\LTcaptype{none} % do not increment counter
\begin{longtable}[]{@{}ll@{}}
\toprule\noalign{}
Situation & Summary \\
\midrule\noalign{}
\endhead
\bottomrule\noalign{}
\endlastfoot
Roughly symmetric, no outliers & mean ± SD \\
Skewed or has outliers & median, IQR, range \\
Categorical & counts and proportions \\
Grouped & Table 1 via \texttt{gtsummary::tbl\_summary()} \\
\end{longtable}
}

\subsection{Contingency tables}\label{contingency-tables}

\begin{Shaded}
\begin{Highlighting}[]
\FunctionTok{table}\NormalTok{(df}\SpecialCharTok{$}\NormalTok{x, df}\SpecialCharTok{$}\NormalTok{y)}
\FunctionTok{prop.table}\NormalTok{(}\FunctionTok{table}\NormalTok{(df}\SpecialCharTok{$}\NormalTok{x, df}\SpecialCharTok{$}\NormalTok{y), }\AttributeTok{margin =} \DecValTok{1}\NormalTok{)}
\end{Highlighting}
\end{Shaded}

\subsection{Bayes' theorem}\label{bayes-theorem}

\(P(A \mid B) = \dfrac{P(B \mid A) P(A)}{P(B)}\)

In odds form: \textbf{posterior odds = prior odds × likelihood ratio.}

\subsection{Diagnostic testing}\label{diagnostic-testing}

{\def\LTcaptype{none} % do not increment counter
\begin{longtable}[]{@{}ll@{}}
\toprule\noalign{}
Metric & Formula \\
\midrule\noalign{}
\endhead
\bottomrule\noalign{}
\endlastfoot
Sensitivity & TP / (TP + FN) \\
Specificity & TN / (TN + FP) \\
PPV & TP / (TP + FP) \\
NPV & TN / (TN + FN) \\
LR+ & Se / (1 − Sp) \\
LR− & (1 − Se) / Sp \\
\end{longtable}
}

PPV collapses at low prevalence --- even a 95\%-accurate test is mostly
false positives.

\subsection{Discrete distributions}\label{discrete-distributions}

{\def\LTcaptype{none} % do not increment counter
\begin{longtable}[]{@{}
  >{\raggedright\arraybackslash}p{(\linewidth - 4\tabcolsep) * \real{0.3333}}
  >{\raggedright\arraybackslash}p{(\linewidth - 4\tabcolsep) * \real{0.3333}}
  >{\raggedright\arraybackslash}p{(\linewidth - 4\tabcolsep) * \real{0.3333}}@{}}
\toprule\noalign{}
\begin{minipage}[b]{\linewidth}\raggedright
Distribution
\end{minipage} & \begin{minipage}[b]{\linewidth}\raggedright
Use
\end{minipage} & \begin{minipage}[b]{\linewidth}\raggedright
R
\end{minipage} \\
\midrule\noalign{}
\endhead
\bottomrule\noalign{}
\endlastfoot
Bernoulli(\emph{p}) & single trial & \texttt{rbinom(n,\ 1,\ p)} \\
Binomial(\emph{n, p}) & \# successes &
\texttt{dbinom\ /\ pbinom\ /\ rbinom} \\
Poisson(\emph{λ}) & rare events & \texttt{dpois\ /\ ppois\ /\ rpois} \\
Negative binomial & overdispersed counts & \texttt{MASS::rnegbin},
\texttt{dnbinom} \\
\end{longtable}
}

Binomial → Poisson as \emph{n} grows and \emph{p} shrinks with \emph{np
= λ}.

\subsection{Continuous distributions}\label{continuous-distributions}

{\def\LTcaptype{none} % do not increment counter
\begin{longtable}[]{@{}ll@{}}
\toprule\noalign{}
Distribution & Use \\
\midrule\noalign{}
\endhead
\bottomrule\noalign{}
\endlastfoot
Normal(\emph{μ, σ}) & almost everything via the CLT \\
Student \emph{t}(ν) & inference about means, small \emph{n} \\
χ²(\emph{ν}) & variance, counts, GoF \\
F(\emph{ν₁, ν₂}) & variance ratios, ANOVA \\
Exponential(\emph{λ}) & time between events, memoryless \\
\end{longtable}
}

R uses the \texttt{d\ /\ p\ /\ q\ /\ r} prefix: density, CDF, quantile,
random.

\subsection{Q-Q plot}\label{q-q-plot}

\begin{Shaded}
\begin{Highlighting}[]
\FunctionTok{ggplot}\NormalTok{(df, }\FunctionTok{aes}\NormalTok{(}\AttributeTok{sample =}\NormalTok{ x)) }\SpecialCharTok{+} \FunctionTok{stat\_qq}\NormalTok{() }\SpecialCharTok{+} \FunctionTok{stat\_qq\_line}\NormalTok{()}
\end{Highlighting}
\end{Shaded}

S-shape → heavy tails. U-shape → skew. Plot before the test.

\subsection{Decision rule}\label{decision-rule-1}

\begin{itemize}
\tightlist
\item
  Reporting a mean? First check a histogram.
\item
  Testing a proportion? First sketch a 2×2 table.
\item
  Choosing a distribution? Simulate before believing.
\end{itemize}

\subsection{Common pitfalls}\label{common-pitfalls-1}

\begin{itemize}
\tightlist
\item
  Quoting PPV without disclosing prevalence.
\item
  Assuming normality because \emph{n} \textgreater{} 30.
\item
  Using mean ± SD on skewed data.
\end{itemize}

\subsection{Further reading}\label{further-reading-1}

\begin{itemize}
\tightlist
\item
  Altman, \emph{Practical Statistics for Medical Research}.
\item
  Gelman et al., \emph{Bayesian Data Analysis}, ch.~1.
\end{itemize}

\subsubsection{Sampling, estimation, and one-sample tests\#\#
Populations vs
samples}\label{sampling-estimation-and-one-sample-tests-populations-vs-samples}

\begin{itemize}
\tightlist
\item
  \textbf{Parameter} = truth about the population (\(\mu\), \(\sigma\),
  \(\pi\)).
\item
  \textbf{Estimator} = statistic computed on a sample (\(\bar x\),
  \(s\), \(\hat p\)).
\item
  \textbf{Bias} = \(E[\hat\theta] - \theta\). \textbf{MSE} = bias² +
  variance.
\end{itemize}

\subsection{Central limit theorem}\label{central-limit-theorem}

For iid samples with finite variance, as \(n \to \infty\):
\(\bar X \sim N\!\left(\mu, \sigma / \sqrt n\right)\)

Holds roughly by \(n \approx 30\) for most non-pathological
distributions.

\subsection{Standard error}\label{standard-error}

\(\mathrm{SE}(\bar X) = \sigma / \sqrt n\), estimated by
\(s / \sqrt n\).

\subsection{Bootstrap}\label{bootstrap}

\begin{Shaded}
\begin{Highlighting}[]
\NormalTok{B }\OtherTok{\textless{}{-}} \DecValTok{2000}
\NormalTok{boot\_mean }\OtherTok{\textless{}{-}} \FunctionTok{replicate}\NormalTok{(B, }\FunctionTok{mean}\NormalTok{(}\FunctionTok{sample}\NormalTok{(x, }\AttributeTok{replace =} \ConstantTok{TRUE}\NormalTok{)))}
\FunctionTok{quantile}\NormalTok{(boot\_mean, }\FunctionTok{c}\NormalTok{(}\FloatTok{0.025}\NormalTok{, }\FloatTok{0.975}\NormalTok{))   }\CommentTok{\# 95\% percentile CI}
\FunctionTok{sd}\NormalTok{(boot\_mean)                           }\CommentTok{\# bootstrap SE}
\end{Highlighting}
\end{Shaded}

\subsection{Permutation test}\label{permutation-test}

\begin{Shaded}
\begin{Highlighting}[]
\NormalTok{obs  }\OtherTok{\textless{}{-}} \FunctionTok{mean}\NormalTok{(a) }\SpecialCharTok{{-}} \FunctionTok{mean}\NormalTok{(b)}
\NormalTok{pool }\OtherTok{\textless{}{-}} \FunctionTok{c}\NormalTok{(a, b)}
\NormalTok{null }\OtherTok{\textless{}{-}} \FunctionTok{replicate}\NormalTok{(}\DecValTok{5000}\NormalTok{, \{}
\NormalTok{  idx }\OtherTok{\textless{}{-}} \FunctionTok{sample}\NormalTok{(}\FunctionTok{seq\_along}\NormalTok{(pool), }\FunctionTok{length}\NormalTok{(a))}
  \FunctionTok{mean}\NormalTok{(pool[idx]) }\SpecialCharTok{{-}} \FunctionTok{mean}\NormalTok{(pool[}\SpecialCharTok{{-}}\NormalTok{idx])}
\NormalTok{\})}
\FunctionTok{mean}\NormalTok{(}\FunctionTok{abs}\NormalTok{(null) }\SpecialCharTok{\textgreater{}=} \FunctionTok{abs}\NormalTok{(obs))   }\CommentTok{\# two{-}sided p}
\end{Highlighting}
\end{Shaded}

\subsection{Maximum likelihood}\label{maximum-likelihood}

Given data \(x\) and model \(f(x; \theta)\):

\begin{itemize}
\tightlist
\item
  \(\ell(\theta) = \sum \log f(x_i; \theta)\).
\item
  \(\hat\theta_{\text{MLE}} = \arg\max_\theta \ell(\theta)\).
\item
  Fisher information \(I(\theta)\); asymptotic SE =
  \(1/\sqrt{I(\hat\theta)}\).
\end{itemize}

\subsection{One-sample tests (five-step
template)}\label{one-sample-tests-five-step-template}

\textbf{Hypothesis → Visualise → Assumptions → Conduct → Conclude.}

{\def\LTcaptype{none} % do not increment counter
\begin{longtable}[]{@{}
  >{\raggedright\arraybackslash}p{(\linewidth - 4\tabcolsep) * \real{0.3333}}
  >{\raggedright\arraybackslash}p{(\linewidth - 4\tabcolsep) * \real{0.3333}}
  >{\raggedright\arraybackslash}p{(\linewidth - 4\tabcolsep) * \real{0.3333}}@{}}
\toprule\noalign{}
\begin{minipage}[b]{\linewidth}\raggedright
Test
\end{minipage} & \begin{minipage}[b]{\linewidth}\raggedright
R
\end{minipage} & \begin{minipage}[b]{\linewidth}\raggedright
Assumptions
\end{minipage} \\
\midrule\noalign{}
\endhead
\bottomrule\noalign{}
\endlastfoot
One-sample \emph{t} & \texttt{t.test(x,\ mu\ =\ 0)} & roughly normal or
large \emph{n} \\
One-proportion & \texttt{prop.test(k,\ n,\ p\ =\ 0.5)} or
\texttt{binom.test} & \(np, n(1-p) \geq 5\) for \texttt{prop.test} \\
\end{longtable}
}

\subsection{Hypothesis-testing
vocabulary}\label{hypothesis-testing-vocabulary}

\begin{itemize}
\tightlist
\item
  \textbf{Type I}: reject true H₀ (probability \(\alpha\)).
\item
  \textbf{Type II}: accept false H₀ (probability \(\beta\)).
\item
  \textbf{Power} = \(1 - \beta\).
\item
  A \emph{p}-value is \textbf{not} the probability that H₀ is true.
\end{itemize}

\subsection{Decision rule}\label{decision-rule-2}

\begin{itemize}
\tightlist
\item
  Want a CI? Bootstrap first, then ask whether a formula applies.
\item
  Want a test? State H₀ and α in writing before running it.
\item
  Want to know if \emph{n} is big enough? Simulate the CLT for your
  outcome.
\end{itemize}

\subsection{Common pitfalls}\label{common-pitfalls-2}

\begin{itemize}
\tightlist
\item
  Reporting SE when you meant SD (or vice versa).
\item
  Using the 95\% CI to ``prove'' there is no effect.
\item
  Running many tests and quoting the smallest \emph{p}-value.
\end{itemize}

\subsection{Further reading}\label{further-reading-2}

\begin{itemize}
\tightlist
\item
  Harrell, \emph{Biostatistics for Biomedical Research}, ch.~3--5.
\item
  Efron \& Tibshirani, \emph{An Introduction to the Bootstrap}.
\end{itemize}

\subsubsection{Two-group comparisons, effect sizes, and power\#\#
Choosing a
comparison}\label{two-group-comparisons-effect-sizes-and-power-choosing-a-comparison}

{\def\LTcaptype{none} % do not increment counter
\begin{longtable}[]{@{}
  >{\raggedright\arraybackslash}p{(\linewidth - 4\tabcolsep) * \real{0.3333}}
  >{\raggedright\arraybackslash}p{(\linewidth - 4\tabcolsep) * \real{0.3333}}
  >{\raggedright\arraybackslash}p{(\linewidth - 4\tabcolsep) * \real{0.3333}}@{}}
\toprule\noalign{}
\begin{minipage}[b]{\linewidth}\raggedright
Design
\end{minipage} & \begin{minipage}[b]{\linewidth}\raggedright
Continuous outcome
\end{minipage} & \begin{minipage}[b]{\linewidth}\raggedright
Binary outcome
\end{minipage} \\
\midrule\noalign{}
\endhead
\bottomrule\noalign{}
\endlastfoot
Two independent groups & \texttt{t.test(y\ \textasciitilde{}\ g)}
(Welch); Wilcoxon rank-sum & \texttt{prop.test} /
\texttt{fisher.test} \\
Paired / pre-post & \texttt{t.test(pre,\ post,\ paired\ =\ TRUE)};
Wilcoxon signed-rank & McNemar \\
\textgreater{} 2 groups, continuous &
\texttt{aov(y\ \textasciitilde{}\ g)}; Kruskal-Wallis & chi-square \\
\end{longtable}
}

\subsection{Effect sizes}\label{effect-sizes}

{\def\LTcaptype{none} % do not increment counter
\begin{longtable}[]{@{}ll@{}}
\toprule\noalign{}
Statistic & What it reports \\
\midrule\noalign{}
\endhead
\bottomrule\noalign{}
\endlastfoot
Mean difference + 95\% CI & primary for continuous \\
Cohen's \emph{d} & standardised mean difference \\
Hedges' \emph{g} & small-sample correction of \emph{d} \\
Risk ratio (RR) / odds ratio (OR) & two proportions \\
Risk difference & absolute, clinically intuitive \\
\end{longtable}
}

\begin{Shaded}
\begin{Highlighting}[]
\NormalTok{effectsize}\SpecialCharTok{::}\FunctionTok{cohens\_d}\NormalTok{(y }\SpecialCharTok{\textasciitilde{}}\NormalTok{ g)}
\end{Highlighting}
\end{Shaded}

\subsection{Two proportions}\label{two-proportions}

\begin{Shaded}
\begin{Highlighting}[]
\FunctionTok{prop.test}\NormalTok{(}\FunctionTok{c}\NormalTok{(tA, tB), }\FunctionTok{c}\NormalTok{(nA, nB))         }\CommentTok{\# asymptotic}
\FunctionTok{fisher.test}\NormalTok{(}\FunctionTok{matrix}\NormalTok{(}\FunctionTok{c}\NormalTok{(tA, nA }\SpecialCharTok{{-}}\NormalTok{ tA,}
\NormalTok{                     tB, nB }\SpecialCharTok{{-}}\NormalTok{ tB), }\DecValTok{2}\NormalTok{))  }\CommentTok{\# exact, small cells}
\end{Highlighting}
\end{Shaded}

Report RR (or OR) with 95\% CI, not just the \emph{p}-value.

\subsection{Correlation}\label{correlation}

{\def\LTcaptype{none} % do not increment counter
\begin{longtable}[]{@{}
  >{\raggedright\arraybackslash}p{(\linewidth - 4\tabcolsep) * \real{0.3333}}
  >{\raggedright\arraybackslash}p{(\linewidth - 4\tabcolsep) * \real{0.3333}}
  >{\raggedright\arraybackslash}p{(\linewidth - 4\tabcolsep) * \real{0.3333}}@{}}
\toprule\noalign{}
\begin{minipage}[b]{\linewidth}\raggedright
Method
\end{minipage} & \begin{minipage}[b]{\linewidth}\raggedright
Captures
\end{minipage} & \begin{minipage}[b]{\linewidth}\raggedright
Assumptions
\end{minipage} \\
\midrule\noalign{}
\endhead
\bottomrule\noalign{}
\endlastfoot
Pearson & linear association, continuous & bivariate normal, no
outliers \\
Spearman & monotonic association & rank-based, robust \\
Kendall & concordant/discordant pairs & robust, slow on large data \\
\end{longtable}
}

\begin{Shaded}
\begin{Highlighting}[]
\FunctionTok{cor.test}\NormalTok{(x, y, }\AttributeTok{method =} \StringTok{"spearman"}\NormalTok{)}
\end{Highlighting}
\end{Shaded}

\subsection{Non-parametric tests}\label{non-parametric-tests}

{\def\LTcaptype{none} % do not increment counter
\begin{longtable}[]{@{}ll@{}}
\toprule\noalign{}
Test & Replaces \\
\midrule\noalign{}
\endhead
\bottomrule\noalign{}
\endlastfoot
Wilcoxon rank-sum (Mann-Whitney) & two-sample \emph{t} \\
Wilcoxon signed-rank & paired \emph{t} \\
Kruskal-Wallis & one-way ANOVA \\
Sign test & paired \emph{t}, when even ranks fail \\
\end{longtable}
}

\subsection{Power and sample size}\label{power-and-sample-size}

\begin{Shaded}
\begin{Highlighting}[]
\NormalTok{pwr}\SpecialCharTok{::}\FunctionTok{pwr.t.test}\NormalTok{(}\AttributeTok{d =} \FloatTok{0.5}\NormalTok{, }\AttributeTok{power =} \FloatTok{0.80}\NormalTok{, }\AttributeTok{sig.level =} \FloatTok{0.05}\NormalTok{,}
                \AttributeTok{type =} \StringTok{"two.sample"}\NormalTok{)}
\end{Highlighting}
\end{Shaded}

{\def\LTcaptype{none} % do not increment counter
\begin{longtable}[]{@{}ll@{}}
\toprule\noalign{}
Family & Function \\
\midrule\noalign{}
\endhead
\bottomrule\noalign{}
\endlastfoot
two-sample \emph{t} & \texttt{pwr.t.test} \\
two proportions & \texttt{pwr.2p.test}, \texttt{pwr.2p2n.test} \\
correlation & \texttt{pwr.r.test} \\
one-way ANOVA & \texttt{pwr.anova.test} \\
\end{longtable}
}

Simulation-based power for anything the textbooks skip:
\texttt{simr::powerSim}.

\subsection{\texorpdfstring{Reporting with
\texttt{gtsummary}}{Reporting with gtsummary}}\label{reporting-with-gtsummary}

\begin{Shaded}
\begin{Highlighting}[]
\FunctionTok{library}\NormalTok{(gtsummary)}
\NormalTok{trial }\SpecialCharTok{|\textgreater{}}
  \FunctionTok{tbl\_summary}\NormalTok{(}\AttributeTok{by =}\NormalTok{ arm, }\AttributeTok{statistic =} \FunctionTok{list}\NormalTok{(}\FunctionTok{all\_continuous}\NormalTok{() }\SpecialCharTok{\textasciitilde{}} \StringTok{"\{mean\} (\{sd\})"}\NormalTok{)) }\SpecialCharTok{|\textgreater{}}
  \FunctionTok{add\_p}\NormalTok{() }\SpecialCharTok{|\textgreater{}}
  \FunctionTok{add\_overall}\NormalTok{()}
\end{Highlighting}
\end{Shaded}

\subsection{Decision rule}\label{decision-rule-3}

\begin{itemize}
\tightlist
\item
  Ask: one-group, two-group, paired, or many-group?
\item
  Check: normal-ish, or should I use a rank-based test?
\item
  Report: point estimate + 95\% CI + effect size; \emph{p}-value last,
  not first.
\end{itemize}

\subsection{Common pitfalls}\label{common-pitfalls-3}

\begin{itemize}
\tightlist
\item
  Equal-variance \emph{t}-test (default in some languages) when
  variances differ.
\item
  Reporting OR when the audience expects RR (or vice versa).
\item
  Forgetting that the \emph{p}-value of a paired test depends on the
  pairing.
\item
  Chaining correlations until one is ``significant''.
\end{itemize}

\subsection{Further reading}\label{further-reading-3}

\begin{itemize}
\tightlist
\item
  Altman, Machin et al., \emph{Statistics with Confidence}.
\item
  \href{https://www.danieldsjoberg.com/gtsummary/}{\texttt{gtsummary}
  docs}.
\end{itemize}

\subsection{Regression}\label{regression}

\subsubsection{Linear regression and diagnostics\#\# Correlation vs
regression}\label{linear-regression-and-diagnostics-correlation-vs-regression}

\begin{itemize}
\tightlist
\item
  Correlation: symmetric; both variables random (\textbf{Model II}).
\item
  Regression: asymmetric; predictor fixed, outcome random (\textbf{Model
  I}).
\item
  Use Model II (e.g.~\texttt{smatr::sma}) when both variables have
  error.
\end{itemize}

\subsection{Simple linear regression}\label{simple-linear-regression}

\(y_i = \beta_0 + \beta_1 x_i + \varepsilon_i, \quad \varepsilon_i \sim N(0, \sigma^2)\)

\begin{Shaded}
\begin{Highlighting}[]
\NormalTok{fit }\OtherTok{\textless{}{-}} \FunctionTok{lm}\NormalTok{(y }\SpecialCharTok{\textasciitilde{}}\NormalTok{ x, }\AttributeTok{data =}\NormalTok{ df)}
\NormalTok{broom}\SpecialCharTok{::}\FunctionTok{tidy}\NormalTok{(fit, }\AttributeTok{conf.int =} \ConstantTok{TRUE}\NormalTok{)}
\NormalTok{broom}\SpecialCharTok{::}\FunctionTok{glance}\NormalTok{(fit)}
\end{Highlighting}
\end{Shaded}

\subsection{Multiple regression}\label{multiple-regression}

\begin{Shaded}
\begin{Highlighting}[]
\NormalTok{fit }\OtherTok{\textless{}{-}} \FunctionTok{lm}\NormalTok{(y }\SpecialCharTok{\textasciitilde{}}\NormalTok{ x1 }\SpecialCharTok{+}\NormalTok{ x2 }\SpecialCharTok{+}\NormalTok{ x1}\SpecialCharTok{:}\NormalTok{x2, }\AttributeTok{data =}\NormalTok{ df)}
\end{Highlighting}
\end{Shaded}

{\def\LTcaptype{none} % do not increment counter
\begin{longtable}[]{@{}ll@{}}
\toprule\noalign{}
Concept & Fix \\
\midrule\noalign{}
\endhead
\bottomrule\noalign{}
\endlastfoot
Confounding & adjust by including the confounder \\
Interaction & include \texttt{x1:x2} (or \texttt{x1\ *\ x2}) \\
Non-linearity & \texttt{splines::ns(x,\ 4)} or \texttt{I(x\^{}2)} \\
Scale sensitivity & centre / standardise predictors \\
\end{longtable}
}

\subsection{Diagnostics}\label{diagnostics}

\begin{Shaded}
\begin{Highlighting}[]
\FunctionTok{plot}\NormalTok{(fit)                   }\CommentTok{\# base: 4 diagnostic plots}
\NormalTok{performance}\SpecialCharTok{::}\FunctionTok{check\_model}\NormalTok{(fit)   }\CommentTok{\# tidy overview}
\NormalTok{car}\SpecialCharTok{::}\FunctionTok{vif}\NormalTok{(fit)               }\CommentTok{\# multicollinearity}
\end{Highlighting}
\end{Shaded}

{\def\LTcaptype{none} % do not increment counter
\begin{longtable}[]{@{}ll@{}}
\toprule\noalign{}
Plot & What it flags \\
\midrule\noalign{}
\endhead
\bottomrule\noalign{}
\endlastfoot
Residuals vs fitted & non-linearity, non-constant variance \\
Normal Q-Q & non-normal residuals \\
Scale-location & heteroscedasticity \\
Leverage / Cook's & influential outliers \\
\end{longtable}
}

VIF \textgreater{} 5 → investigate collinearity; \textgreater{} 10 → fix
it.

\subsection{Robust / weighted}\label{robust-weighted}

\begin{Shaded}
\begin{Highlighting}[]
\NormalTok{MASS}\SpecialCharTok{::}\FunctionTok{rlm}\NormalTok{(y }\SpecialCharTok{\textasciitilde{}}\NormalTok{ x, }\AttributeTok{data =}\NormalTok{ df)          }\CommentTok{\# robust to outliers}
\FunctionTok{lm}\NormalTok{(y }\SpecialCharTok{\textasciitilde{}}\NormalTok{ x, }\AttributeTok{weights =} \DecValTok{1} \SpecialCharTok{/}\NormalTok{ var\_i)        }\CommentTok{\# weighted LS}

\CommentTok{\# Heteroscedasticity{-}consistent SEs}
\FunctionTok{library}\NormalTok{(sandwich); }\FunctionTok{library}\NormalTok{(lmtest)}
\FunctionTok{coeftest}\NormalTok{(fit, }\AttributeTok{vcov =} \FunctionTok{vcovHC}\NormalTok{(fit, }\AttributeTok{type =} \StringTok{"HC3"}\NormalTok{))}
\end{Highlighting}
\end{Shaded}

\subsection{Decision rule}\label{decision-rule-4}

\begin{itemize}
\tightlist
\item
  Residual plot ugly? Transform, add terms, or switch link.
\item
  VIF explodes? Drop one of the collinear predictors or combine them.
\item
  Outlier with high Cook's distance? Investigate before deleting.
\item
  Variance depends on X? HC3 SEs or weighted LS.
\end{itemize}

\subsection{Common pitfalls}\label{common-pitfalls-4}

\begin{itemize}
\tightlist
\item
  Interpreting \(\beta_1\) as ``effect of X'' under confounding.
\item
  Reporting \(R^2\) on in-sample data as if it were generalisable.
\item
  Including an interaction but only reporting main effects.
\item
  Centring categorical predictors (don't --- use reference coding).
\end{itemize}

\subsection{Further reading}\label{further-reading-4}

\begin{itemize}
\tightlist
\item
  Harrell, \emph{Regression Modeling Strategies}, ch.~4--5.
\item
  Fox, \emph{Applied Regression Analysis}.
\end{itemize}

\subsubsection{ANOVA, mixed models, and non-linear regression\#\#
One-way
ANOVA}\label{anova-mixed-models-and-non-linear-regression-one-way-anova}

\begin{Shaded}
\begin{Highlighting}[]
\FunctionTok{aov}\NormalTok{(y }\SpecialCharTok{\textasciitilde{}}\NormalTok{ group, }\AttributeTok{data =}\NormalTok{ df) }\SpecialCharTok{|\textgreater{}} \FunctionTok{summary}\NormalTok{()}
\end{Highlighting}
\end{Shaded}

ANOVA is a linear model with a categorical predictor. The F-test
compares between-group to within-group variance.

\subsection{\texorpdfstring{Contrasts with
\texttt{emmeans}}{Contrasts with emmeans}}\label{contrasts-with-emmeans}

\begin{Shaded}
\begin{Highlighting}[]
\FunctionTok{library}\NormalTok{(emmeans)}
\NormalTok{emm }\OtherTok{\textless{}{-}} \FunctionTok{emmeans}\NormalTok{(fit, }\SpecialCharTok{\textasciitilde{}}\NormalTok{ group)}
\FunctionTok{pairs}\NormalTok{(emm, }\AttributeTok{adjust =} \StringTok{"tukey"}\NormalTok{)   }\CommentTok{\# all pairwise}
\FunctionTok{contrast}\NormalTok{(emm, }\FunctionTok{list}\NormalTok{(}\AttributeTok{TrtVsCtrl =} \FunctionTok{c}\NormalTok{(}\SpecialCharTok{{-}}\DecValTok{1}\NormalTok{, }\DecValTok{1}\NormalTok{, }\DecValTok{1}\NormalTok{, }\DecValTok{1}\NormalTok{) }\SpecialCharTok{/} \DecValTok{3}\NormalTok{))}
\end{Highlighting}
\end{Shaded}

Pre-specify contrasts before looking at the data; correct for
multiplicity.

\subsection{Two-way / factorial ANOVA}\label{two-way-factorial-anova}

\begin{Shaded}
\begin{Highlighting}[]
\FunctionTok{aov}\NormalTok{(y }\SpecialCharTok{\textasciitilde{}}\NormalTok{ A }\SpecialCharTok{*}\NormalTok{ B, }\AttributeTok{data =}\NormalTok{ df) }\SpecialCharTok{|\textgreater{}} \FunctionTok{summary}\NormalTok{()}
\FunctionTok{emmip}\NormalTok{(fit, A }\SpecialCharTok{\textasciitilde{}}\NormalTok{ B)   }\CommentTok{\# interaction plot}
\end{Highlighting}
\end{Shaded}

Interaction means ``effect of A differs by level of B''. Report the
interaction first; main effects are conditional.

\subsection{Repeated measures /
blocking}\label{repeated-measures-blocking}

\begin{itemize}
\tightlist
\item
  RCBD: \texttt{aov(y\ \textasciitilde{}\ treatment\ +\ Error(block))}.
\item
  Repeated measures: move to a mixed model.
\end{itemize}

\begin{Shaded}
\begin{Highlighting}[]
\FunctionTok{library}\NormalTok{(lme4); }\FunctionTok{library}\NormalTok{(lmerTest)}
\FunctionTok{lmer}\NormalTok{(y }\SpecialCharTok{\textasciitilde{}}\NormalTok{ treatment }\SpecialCharTok{+}\NormalTok{ time }\SpecialCharTok{+}\NormalTok{ (}\DecValTok{1} \SpecialCharTok{|}\NormalTok{ subject), }\AttributeTok{data =}\NormalTok{ df)}
\end{Highlighting}
\end{Shaded}

\subsection{GAMs --- smooth non-linear
terms}\label{gams-smooth-non-linear-terms}

\begin{Shaded}
\begin{Highlighting}[]
\FunctionTok{library}\NormalTok{(mgcv); }\FunctionTok{library}\NormalTok{(gratia)}
\NormalTok{fit }\OtherTok{\textless{}{-}} \FunctionTok{gam}\NormalTok{(y }\SpecialCharTok{\textasciitilde{}} \FunctionTok{s}\NormalTok{(x, }\AttributeTok{k =} \DecValTok{10}\NormalTok{) }\SpecialCharTok{+}\NormalTok{ z, }\AttributeTok{data =}\NormalTok{ df)}
\FunctionTok{summary}\NormalTok{(fit)      }\CommentTok{\# edf tells you how "wiggly"}
\FunctionTok{draw}\NormalTok{(fit)         }\CommentTok{\# smooth + CI}
\end{Highlighting}
\end{Shaded}

edf ≈ 1 → nearly linear; \textgreater{} 4 → clearly non-linear.

\subsection{\texorpdfstring{Non-linear regression
(\texttt{nls})}{Non-linear regression (nls)}}\label{non-linear-regression-nls}

\begin{Shaded}
\begin{Highlighting}[]
\CommentTok{\# Michaelis{-}Menten: y = Vmax * x / (K + x)}
\NormalTok{fit }\OtherTok{\textless{}{-}} \FunctionTok{nls}\NormalTok{(y }\SpecialCharTok{\textasciitilde{}}\NormalTok{ Vmax }\SpecialCharTok{*}\NormalTok{ x }\SpecialCharTok{/}\NormalTok{ (K }\SpecialCharTok{+}\NormalTok{ x),}
           \AttributeTok{data =}\NormalTok{ df, }\AttributeTok{start =} \FunctionTok{list}\NormalTok{(}\AttributeTok{Vmax =} \DecValTok{1}\NormalTok{, }\AttributeTok{K =} \DecValTok{1}\NormalTok{))}
\end{Highlighting}
\end{Shaded}

Start values matter. If it fails, plot first to guess reasonable starts.

\subsection{Decision rule}\label{decision-rule-5}

\begin{itemize}
\tightlist
\item
  Categorical predictor, \textgreater{} 2 levels → ANOVA + contrasts.
\item
  Factorial design → include interaction, report it first.
\item
  Effect obviously curved → GAM with spline; else try \texttt{nls}.
\item
  Repeated measures → mixed model, not repeated-measures ANOVA.
\end{itemize}

\subsection{Common pitfalls}\label{common-pitfalls-5}

\begin{itemize}
\tightlist
\item
  Tukey HSD without pre-specified contrasts of interest.
\item
  ANOVA p \textless{} 0.05 reported alone --- without naming
  \emph{which} groups differ.
\item
  Forcing a GAM onto monotonic data that \texttt{nls} fits cleanly.
\item
  Ignoring the random effect in clustered designs (pseudoreplication).
\end{itemize}

\subsection{Further reading}\label{further-reading-5}

\begin{itemize}
\tightlist
\item
  Wood, \emph{Generalized Additive Models}, 2e.
\item
  \href{https://cran.r-project.org/package=emmeans}{\texttt{emmeans}
  vignette}.
\end{itemize}

\subsubsection{Generalised linear models and prediction evaluation\#\#
GLM link
functions}\label{generalised-linear-models-and-prediction-evaluation-glm-link-functions}

{\def\LTcaptype{none} % do not increment counter
\begin{longtable}[]{@{}
  >{\raggedright\arraybackslash}p{(\linewidth - 6\tabcolsep) * \real{0.2500}}
  >{\raggedright\arraybackslash}p{(\linewidth - 6\tabcolsep) * \real{0.2500}}
  >{\raggedright\arraybackslash}p{(\linewidth - 6\tabcolsep) * \real{0.2500}}
  >{\raggedright\arraybackslash}p{(\linewidth - 6\tabcolsep) * \real{0.2500}}@{}}
\toprule\noalign{}
\begin{minipage}[b]{\linewidth}\raggedright
Outcome
\end{minipage} & \begin{minipage}[b]{\linewidth}\raggedright
Family
\end{minipage} & \begin{minipage}[b]{\linewidth}\raggedright
Link
\end{minipage} & \begin{minipage}[b]{\linewidth}\raggedright
R
\end{minipage} \\
\midrule\noalign{}
\endhead
\bottomrule\noalign{}
\endlastfoot
Continuous & gaussian & identity & \texttt{lm()} \\
Binary & binomial & logit &
\texttt{glm(y\ \textasciitilde{}\ x,\ family\ =\ binomial)} \\
Count & poisson & log &
\texttt{glm(y\ \textasciitilde{}\ x,\ family\ =\ poisson,\ offset\ =\ log(t))} \\
Overdispersed count & neg. binomial & log & \texttt{MASS::glm.nb} \\
Ordinal & cumulative & logit & \texttt{MASS::polr} \\
Nominal & multinomial & logit & \texttt{nnet::multinom} \\
\end{longtable}
}

\subsection{Logistic regression}\label{logistic-regression}

\begin{Shaded}
\begin{Highlighting}[]
\NormalTok{fit }\OtherTok{\textless{}{-}} \FunctionTok{glm}\NormalTok{(y }\SpecialCharTok{\textasciitilde{}}\NormalTok{ x1 }\SpecialCharTok{+}\NormalTok{ x2, }\AttributeTok{data =}\NormalTok{ df, }\AttributeTok{family =}\NormalTok{ binomial)}
\FunctionTok{exp}\NormalTok{(}\FunctionTok{coef}\NormalTok{(fit))                        }\CommentTok{\# odds ratios}
\FunctionTok{exp}\NormalTok{(}\FunctionTok{confint}\NormalTok{(fit))                     }\CommentTok{\# 95\% CI on OR}
\FunctionTok{predict}\NormalTok{(fit, }\AttributeTok{newdata =}\NormalTok{ nd, }\AttributeTok{type =} \StringTok{"response"}\NormalTok{)}
\end{Highlighting}
\end{Shaded}

\begin{itemize}
\tightlist
\item
  Check for perfect separation (huge SEs).
\item
  Interpret ORs cautiously; RR is more intuitive for the audience.
\end{itemize}

\subsection{ANCOVA in an RCT}\label{ancova-in-an-rct}

Adjust for baseline; \textbf{do not} analyse the change score.

\begin{Shaded}
\begin{Highlighting}[]
\FunctionTok{lm}\NormalTok{(y\_followup }\SpecialCharTok{\textasciitilde{}}\NormalTok{ arm }\SpecialCharTok{+}\NormalTok{ y\_baseline, }\AttributeTok{data =}\NormalTok{ trial)}
\end{Highlighting}
\end{Shaded}

More efficient than simple t-test on change scores when baseline and
follow-up are correlated.

\subsection{Poisson / negative
binomial}\label{poisson-negative-binomial}

\begin{Shaded}
\begin{Highlighting}[]
\FunctionTok{glm}\NormalTok{(cases }\SpecialCharTok{\textasciitilde{}}\NormalTok{ x }\SpecialCharTok{+} \FunctionTok{offset}\NormalTok{(}\FunctionTok{log}\NormalTok{(person\_years)),}
    \AttributeTok{family =}\NormalTok{ poisson, }\AttributeTok{data =}\NormalTok{ df)}
\NormalTok{MASS}\SpecialCharTok{::}\FunctionTok{glm.nb}\NormalTok{(cases }\SpecialCharTok{\textasciitilde{}}\NormalTok{ x }\SpecialCharTok{+} \FunctionTok{offset}\NormalTok{(}\FunctionTok{log}\NormalTok{(person\_years)), }\AttributeTok{data =}\NormalTok{ df)}
\end{Highlighting}
\end{Shaded}

Check dispersion:
\texttt{sum(residuals(fit,\ type\ =\ "pearson")\^{}2)\ /\ df.residual}.
If \textgreater{} 1.5, switch to NB.

\subsection{Evaluation --- calibration +
discrimination}\label{evaluation-calibration-discrimination}

{\def\LTcaptype{none} % do not increment counter
\begin{longtable}[]{@{}
  >{\raggedright\arraybackslash}p{(\linewidth - 4\tabcolsep) * \real{0.3333}}
  >{\raggedright\arraybackslash}p{(\linewidth - 4\tabcolsep) * \real{0.3333}}
  >{\raggedright\arraybackslash}p{(\linewidth - 4\tabcolsep) * \real{0.3333}}@{}}
\toprule\noalign{}
\begin{minipage}[b]{\linewidth}\raggedright
Metric
\end{minipage} & \begin{minipage}[b]{\linewidth}\raggedright
Means
\end{minipage} & \begin{minipage}[b]{\linewidth}\raggedright
R
\end{minipage} \\
\midrule\noalign{}
\endhead
\bottomrule\noalign{}
\endlastfoot
Calibration plot & predicted vs observed & \texttt{rms::val.prob},
manual bin \\
ROC / AUC & rank ordering & \texttt{pROC::roc(y,\ phat)} \\
Brier score & overall accuracy & \texttt{mean((phat\ -\ y)\^{}2)} \\
Calibration slope / intercept & systematic bias & from logistic
recalibration \\
\end{longtable}
}

\begin{Shaded}
\begin{Highlighting}[]
\FunctionTok{library}\NormalTok{(pROC)}
\NormalTok{roc\_obj }\OtherTok{\textless{}{-}} \FunctionTok{roc}\NormalTok{(y, phat)}
\FunctionTok{auc}\NormalTok{(roc\_obj); }\FunctionTok{ci.auc}\NormalTok{(roc\_obj)}
\FunctionTok{plot}\NormalTok{(roc\_obj)}
\end{Highlighting}
\end{Shaded}

\subsection{Decision rule}\label{decision-rule-6}

\begin{itemize}
\tightlist
\item
  Binary outcome → logistic; report OR + 95\% CI.
\item
  Count outcome → Poisson; check overdispersion; NB if needed.
\item
  Trial analysis → ANCOVA, not change score.
\item
  Prediction model → calibration curve first, ROC second, decision
  curves third.
\end{itemize}

\subsection{Common pitfalls}\label{common-pitfalls-6}

\begin{itemize}
\tightlist
\item
  Quoting AUC without calibration (a discriminating but miscalibrated
  model is dangerous).
\item
  Ignoring offsets in count data.
\item
  Using ordinal logit when the proportional-odds assumption fails.
\item
  Presenting logistic regression coefficients on the log-odds scale
  without OR.
\end{itemize}

\subsection{Further reading}\label{further-reading-6}

\begin{itemize}
\tightlist
\item
  Harrell, \emph{RMS}, ch.~10--12.
\item
  Steyerberg, \emph{Clinical Prediction Models}.
\end{itemize}

\subsubsection{Agreement, survival primer, and reporting\#\# Don't
dichotomise}\label{agreement-survival-primer-and-reporting-dont-dichotomise}

Cutting a continuous variable at the median loses ≈ 37\% of the
information for a linear association. Keep predictors continuous; model
non-linearity with a spline if needed.

\subsection{Change scores vs ANCOVA}\label{change-scores-vs-ancova}

\begin{itemize}
\tightlist
\item
  Change scores suffer from \textbf{regression to the mean}.
\item
  ANCOVA (adjust for baseline) is the efficient analysis in an RCT.
\item
  In observational studies with unequal baselines, both approaches
  answer subtly different questions --- state which.
\end{itemize}

\subsection{Agreement \& reliability}\label{agreement-reliability}

{\def\LTcaptype{none} % do not increment counter
\begin{longtable}[]{@{}
  >{\raggedright\arraybackslash}p{(\linewidth - 4\tabcolsep) * \real{0.3333}}
  >{\raggedright\arraybackslash}p{(\linewidth - 4\tabcolsep) * \real{0.3333}}
  >{\raggedright\arraybackslash}p{(\linewidth - 4\tabcolsep) * \real{0.3333}}@{}}
\toprule\noalign{}
\begin{minipage}[b]{\linewidth}\raggedright
Statistic
\end{minipage} & \begin{minipage}[b]{\linewidth}\raggedright
Use
\end{minipage} & \begin{minipage}[b]{\linewidth}\raggedright
R
\end{minipage} \\
\midrule\noalign{}
\endhead
\bottomrule\noalign{}
\endlastfoot
Cohen's κ & two raters, categorical &
\texttt{irr::kappa2(cbind(r1,\ r2))} \\
Weighted κ & ordinal categorical &
\texttt{irr::kappa2(...,\ weight\ =\ "squared")} \\
ICC(2,1) / ICC(3,1) & continuous & \texttt{psych::ICC(df)} \\
Bland-Altman & two methods, continuous & plot \(\bar{xy}\) vs
\(y - x\) \\
\end{longtable}
}

\begin{Shaded}
\begin{Highlighting}[]
\NormalTok{mean\_diff }\OtherTok{\textless{}{-}} \FunctionTok{mean}\NormalTok{(y }\SpecialCharTok{{-}}\NormalTok{ x)}
\NormalTok{loa }\OtherTok{\textless{}{-}}\NormalTok{ mean\_diff }\SpecialCharTok{+} \FunctionTok{c}\NormalTok{(}\SpecialCharTok{{-}}\FloatTok{1.96}\NormalTok{, }\FloatTok{1.96}\NormalTok{) }\SpecialCharTok{*} \FunctionTok{sd}\NormalTok{(y }\SpecialCharTok{{-}}\NormalTok{ x)}
\FunctionTok{ggplot}\NormalTok{(df, }\FunctionTok{aes}\NormalTok{((x }\SpecialCharTok{+}\NormalTok{ y)}\SpecialCharTok{/}\DecValTok{2}\NormalTok{, y }\SpecialCharTok{{-}}\NormalTok{ x)) }\SpecialCharTok{+} \FunctionTok{geom\_point}\NormalTok{() }\SpecialCharTok{+}
  \FunctionTok{geom\_hline}\NormalTok{(}\AttributeTok{yintercept =} \FunctionTok{c}\NormalTok{(mean\_diff, loa), }\AttributeTok{linetype =} \DecValTok{2}\NormalTok{)}
\end{Highlighting}
\end{Shaded}

\subsection{Survival primer}\label{survival-primer}

{\def\LTcaptype{none} % do not increment counter
\begin{longtable}[]{@{}
  >{\raggedright\arraybackslash}p{(\linewidth - 2\tabcolsep) * \real{0.5000}}
  >{\raggedright\arraybackslash}p{(\linewidth - 2\tabcolsep) * \real{0.5000}}@{}}
\toprule\noalign{}
\begin{minipage}[b]{\linewidth}\raggedright
Concept
\end{minipage} & \begin{minipage}[b]{\linewidth}\raggedright
R
\end{minipage} \\
\midrule\noalign{}
\endhead
\bottomrule\noalign{}
\endlastfoot
Kaplan-Meier &
\texttt{survfit(Surv(time,\ event)\ \textasciitilde{}\ g)} +
\texttt{ggsurvfit::ggsurvfit} \\
Log-rank test &
\texttt{survdiff(Surv(time,\ event)\ \textasciitilde{}\ g)} \\
Cox PH model &
\texttt{coxph(Surv(time,\ event)\ \textasciitilde{}\ x1\ +\ x2)} \\
PH check & \texttt{cox.zph(fit)} \\
\end{longtable}
}

\begin{itemize}
\tightlist
\item
  Report a hazard ratio with 95\% CI.
\item
  PH violated? → time-varying coefficient or stratify.
\item
  Interpret HR only after checking for non-informative censoring.
\end{itemize}

\subsection{Decision curves, NRI, IDI}\label{decision-curves-nri-idi}

\begin{itemize}
\tightlist
\item
  \textbf{Decision curve}: net benefit vs threshold probability ---
  dominates ``treat all'' and ``treat none'' when useful.
\item
  \textbf{NRI}: how many events move up / non-events move down in risk.
\item
  \textbf{IDI}: mean change in predicted probability by class.
\item
  Report decision curve first; NRI/IDI as secondary.
\end{itemize}

\subsection{Explanation vs prediction
(Shmueli)}\label{explanation-vs-prediction-shmueli}

{\def\LTcaptype{none} % do not increment counter
\begin{longtable}[]{@{}
  >{\raggedright\arraybackslash}p{(\linewidth - 2\tabcolsep) * \real{0.5000}}
  >{\raggedright\arraybackslash}p{(\linewidth - 2\tabcolsep) * \real{0.5000}}@{}}
\toprule\noalign{}
\begin{minipage}[b]{\linewidth}\raggedright
Explanatory
\end{minipage} & \begin{minipage}[b]{\linewidth}\raggedright
Predictive
\end{minipage} \\
\midrule\noalign{}
\endhead
\bottomrule\noalign{}
\endlastfoot
Goal: inference about β & Goal: minimise out-of-sample loss \\
Tools: ANOVA, diagnostics, intervals & Tools: CV, regularisation,
ensembles \\
Metric: interval coverage, bias & Metric: RMSE, AUC, Brier on
hold-out \\
\end{longtable}
}

\subsection{Reporting guidelines}\label{reporting-guidelines}

{\def\LTcaptype{none} % do not increment counter
\begin{longtable}[]{@{}ll@{}}
\toprule\noalign{}
Design & Guideline \\
\midrule\noalign{}
\endhead
\bottomrule\noalign{}
\endlastfoot
Randomised trial & CONSORT \\
Observational & STROBE \\
Diagnostic-accuracy & STARD \\
Prediction model (incl.~AI) & TRIPOD / TRIPOD-AI \\
Systematic review & PRISMA \\
Animal research & ARRIVE \\
\end{longtable}
}

\subsection{Decision rule}\label{decision-rule-7}

\begin{itemize}
\tightlist
\item
  Continuous predictor + continuous outcome → no median splits.
\item
  Two-group with repeated measurement → ANCOVA.
\item
  New rater / device → Bland-Altman + ICC.
\item
  Time-to-event → KM + Cox, always check PH.
\item
  Prediction model → TRIPOD checklist at submission.
\end{itemize}

\subsection{Common pitfalls}\label{common-pitfalls-7}

\begin{itemize}
\tightlist
\item
  Reporting κ on near-constant outcomes (high \% agreement, low κ).
\item
  Quoting median survival that is never reached.
\item
  Forgetting to mention PH assumption checking.
\item
  Claiming a ``prediction model'' built on the same data used for
  evaluation.
\end{itemize}

\subsection{Further reading}\label{further-reading-7}

\begin{itemize}
\tightlist
\item
  Harrell, \emph{BBR}, ch.~17--18.
\item
  Royston \& Altman, \emph{Prognosis research}.
\end{itemize}

\subsection{Design \& Causal Inference}\label{design-causal-inference}

\subsubsection{Study designs and power\#\# Observational
designs}\label{study-designs-and-power-observational-designs}

{\def\LTcaptype{none} % do not increment counter
\begin{longtable}[]{@{}
  >{\raggedright\arraybackslash}p{(\linewidth - 4\tabcolsep) * \real{0.3333}}
  >{\raggedright\arraybackslash}p{(\linewidth - 4\tabcolsep) * \real{0.3333}}
  >{\raggedright\arraybackslash}p{(\linewidth - 4\tabcolsep) * \real{0.3333}}@{}}
\toprule\noalign{}
\begin{minipage}[b]{\linewidth}\raggedright
Design
\end{minipage} & \begin{minipage}[b]{\linewidth}\raggedright
Strengths
\end{minipage} & \begin{minipage}[b]{\linewidth}\raggedright
Weaknesses
\end{minipage} \\
\midrule\noalign{}
\endhead
\bottomrule\noalign{}
\endlastfoot
Cohort & temporality, incidence & expensive, loss to follow-up \\
Case-control & rare outcomes, cheap & recall + selection bias \\
Cross-sectional & prevalence, screening & no temporality \\
Case-crossover & within-person comparison & for transient exposures
only \\
\end{longtable}
}

STROBE checklist: \url{https://www.strobe-statement.org/}

\subsection{Trial designs}\label{trial-designs}

{\def\LTcaptype{none} % do not increment counter
\begin{longtable}[]{@{}
  >{\raggedright\arraybackslash}p{(\linewidth - 2\tabcolsep) * \real{0.5000}}
  >{\raggedright\arraybackslash}p{(\linewidth - 2\tabcolsep) * \real{0.5000}}@{}}
\toprule\noalign{}
\begin{minipage}[b]{\linewidth}\raggedright
Type
\end{minipage} & \begin{minipage}[b]{\linewidth}\raggedright
When
\end{minipage} \\
\midrule\noalign{}
\endhead
\bottomrule\noalign{}
\endlastfoot
Parallel-group & most common; independent arms \\
Crossover & stable chronic condition, carry-over washable \\
Cluster & intervention at group level (schools, clinics) \\
Factorial & two or more interventions, test interactions \\
Adaptive & pre-specified modifications based on interim data \\
Non-inferiority & new treatment not worse than control by margin
\(\Delta\) \\
Equivalence & two-sided non-inferiority \\
\end{longtable}
}

CONSORT for RCTs: \url{https://www.consort-statement.org/}

\subsection{Bench / translational}\label{bench-translational}

\begin{itemize}
\tightlist
\item
  \textbf{Blocking} reduces nuisance variation (plate, day, operator).
\item
  \textbf{Factorial} tests interactions efficiently.
\item
  \textbf{Split-plot} handles two levels of randomisation.
\item
  \textbf{Pseudoreplication}: technical replicates ≠ biological
  replicates.
\end{itemize}

\subsection{Power --- closed form}\label{power-closed-form}

\begin{Shaded}
\begin{Highlighting}[]
\FunctionTok{library}\NormalTok{(pwr)}
\FunctionTok{pwr.t.test}\NormalTok{(}\AttributeTok{d =} \FloatTok{0.5}\NormalTok{, }\AttributeTok{power =} \FloatTok{0.80}\NormalTok{, }\AttributeTok{sig.level =} \FloatTok{0.05}\NormalTok{,}
           \AttributeTok{type =} \StringTok{"two.sample"}\NormalTok{)}
\FunctionTok{pwr.2p.test}\NormalTok{(}\AttributeTok{h =} \FunctionTok{ES.h}\NormalTok{(}\FloatTok{0.3}\NormalTok{, }\FloatTok{0.2}\NormalTok{), }\AttributeTok{power =} \FloatTok{0.80}\NormalTok{)}
\FunctionTok{pwr.r.test}\NormalTok{(}\AttributeTok{r =} \FloatTok{0.3}\NormalTok{, }\AttributeTok{power =} \FloatTok{0.80}\NormalTok{)}
\FunctionTok{pwr.anova.test}\NormalTok{(}\AttributeTok{k =} \DecValTok{4}\NormalTok{, }\AttributeTok{f =} \FloatTok{0.25}\NormalTok{, }\AttributeTok{power =} \FloatTok{0.80}\NormalTok{)}
\end{Highlighting}
\end{Shaded}

Effect-size conventions (Cohen): small \(d\) = 0.2, medium 0.5, large
0.8.

\subsection{Power --- simulation}\label{power-simulation}

\begin{Shaded}
\begin{Highlighting}[]
\FunctionTok{library}\NormalTok{(simr)}
\CommentTok{\# Build a pilot model, increase N, or tweak fixef, then:}
\NormalTok{ps }\OtherTok{\textless{}{-}} \FunctionTok{powerSim}\NormalTok{(model, }\AttributeTok{nsim =} \DecValTok{500}\NormalTok{, }\AttributeTok{test =} \FunctionTok{fixed}\NormalTok{(}\StringTok{"arm"}\NormalTok{))}
\NormalTok{ps}
\end{Highlighting}
\end{Shaded}

Simulation wins for any design the textbook skips: mixed models,
adaptive rules, non-standard outcomes.

\subsection{Decision rule}\label{decision-rule-8}

\begin{itemize}
\tightlist
\item
  Randomise if you can. If not, draw a DAG and name the biases.
\item
  Power calculation before the protocol freeze, not after data
  collection.
\item
  Cluster randomisation → design effect \(1 + (\bar m - 1)\rho\);
  inflate N.
\item
  Bench experiments → treat batch, plate, and operator as random
  effects.
\end{itemize}

\subsection{Common pitfalls}\label{common-pitfalls-8}

\begin{itemize}
\tightlist
\item
  Planning a cluster RCT without inflating N for design effect.
\item
  Borrowing a pilot effect size without acknowledging noise.
\item
  Running a non-inferiority trial as if it were superiority.
\item
  Treating technical replicates as if they were biological.
\end{itemize}

\subsection{Further reading}\label{further-reading-8}

\begin{itemize}
\tightlist
\item
  Matthews, \emph{Introduction to Randomized Controlled Clinical
  Trials}.
\item
  Chow \& Liu, \emph{Design and Analysis of Clinical Trials}.
\end{itemize}

\subsubsection{Missing data, mixed models, and longitudinal\#\#
Missingness
mechanisms}\label{missing-data-mixed-models-and-longitudinal-missingness-mechanisms}

{\def\LTcaptype{none} % do not increment counter
\begin{longtable}[]{@{}
  >{\raggedright\arraybackslash}p{(\linewidth - 4\tabcolsep) * \real{0.3333}}
  >{\raggedright\arraybackslash}p{(\linewidth - 4\tabcolsep) * \real{0.3333}}
  >{\raggedright\arraybackslash}p{(\linewidth - 4\tabcolsep) * \real{0.3333}}@{}}
\toprule\noalign{}
\begin{minipage}[b]{\linewidth}\raggedright
Mechanism
\end{minipage} & \begin{minipage}[b]{\linewidth}\raggedright
Definition
\end{minipage} & \begin{minipage}[b]{\linewidth}\raggedright
Complete-case analysis?
\end{minipage} \\
\midrule\noalign{}
\endhead
\bottomrule\noalign{}
\endlastfoot
MCAR & missingness independent of data & unbiased but inefficient \\
MAR & missingness depends on \textbf{observed} data & biased, imputation
helps \\
MNAR & missingness depends on \textbf{unobserved} values & imputation
under model only \\
\end{longtable}
}

Test is impossible without assumptions --- reason about the mechanism.

\subsection{\texorpdfstring{Multiple imputation with
\texttt{mice}}{Multiple imputation with mice}}\label{multiple-imputation-with-mice}

\begin{Shaded}
\begin{Highlighting}[]
\FunctionTok{library}\NormalTok{(mice)}
\NormalTok{imp }\OtherTok{\textless{}{-}} \FunctionTok{mice}\NormalTok{(df, }\AttributeTok{m =} \DecValTok{20}\NormalTok{, }\AttributeTok{seed =} \DecValTok{42}\NormalTok{, }\AttributeTok{printFlag =} \ConstantTok{FALSE}\NormalTok{)}
\NormalTok{fit }\OtherTok{\textless{}{-}} \FunctionTok{with}\NormalTok{(imp, }\FunctionTok{lm}\NormalTok{(y }\SpecialCharTok{\textasciitilde{}}\NormalTok{ x1 }\SpecialCharTok{+}\NormalTok{ x2))}
\NormalTok{pooled }\OtherTok{\textless{}{-}} \FunctionTok{pool}\NormalTok{(fit)}
\FunctionTok{summary}\NormalTok{(pooled, }\AttributeTok{conf.int =} \ConstantTok{TRUE}\NormalTok{)}
\end{Highlighting}
\end{Shaded}

Rule of thumb: \emph{m} ≥ 100 × fraction of missing information, but 20
is a fine starting point.

\subsection{Linear mixed models}\label{linear-mixed-models}

\begin{Shaded}
\begin{Highlighting}[]
\FunctionTok{library}\NormalTok{(lme4); }\FunctionTok{library}\NormalTok{(lmerTest)}
\FunctionTok{lmer}\NormalTok{(y }\SpecialCharTok{\textasciitilde{}}\NormalTok{ treatment }\SpecialCharTok{*}\NormalTok{ time }\SpecialCharTok{+}\NormalTok{ (}\DecValTok{1} \SpecialCharTok{|}\NormalTok{ subject), }\AttributeTok{data =}\NormalTok{ df)}
\FunctionTok{lmer}\NormalTok{(y }\SpecialCharTok{\textasciitilde{}}\NormalTok{ treatment }\SpecialCharTok{*}\NormalTok{ time }\SpecialCharTok{+}\NormalTok{ (time }\SpecialCharTok{|}\NormalTok{ subject), }\AttributeTok{data =}\NormalTok{ df)  }\CommentTok{\# random slope}
\end{Highlighting}
\end{Shaded}

\begin{itemize}
\tightlist
\item
  Random intercept: each cluster has its own baseline.
\item
  Random slope: each cluster has its own trajectory.
\item
  ICC \(\rho = \sigma_u^2 / (\sigma_u^2 + \sigma_e^2)\).
\end{itemize}

REML for variance components; ML only for likelihood-ratio tests on
fixed effects.

\subsection{GLMMs \& GEE}\label{glmms-gee}

\begin{Shaded}
\begin{Highlighting}[]
\FunctionTok{library}\NormalTok{(lme4)}
\FunctionTok{glmer}\NormalTok{(y }\SpecialCharTok{\textasciitilde{}}\NormalTok{ x }\SpecialCharTok{+}\NormalTok{ (}\DecValTok{1} \SpecialCharTok{|}\NormalTok{ cluster), }\AttributeTok{data =}\NormalTok{ df, }\AttributeTok{family =}\NormalTok{ binomial)}

\FunctionTok{library}\NormalTok{(geepack)}
\FunctionTok{geeglm}\NormalTok{(y }\SpecialCharTok{\textasciitilde{}}\NormalTok{ x, }\AttributeTok{id =}\NormalTok{ cluster, }\AttributeTok{data =}\NormalTok{ df,}
       \AttributeTok{family =}\NormalTok{ binomial, }\AttributeTok{corstr =} \StringTok{"exchangeable"}\NormalTok{)}
\end{Highlighting}
\end{Shaded}

\begin{itemize}
\tightlist
\item
  GLMM: subject-specific effects, good for prediction.
\item
  GEE: population-average effects, robust to working-correlation
  misspecification (when combined with robust SEs).
\end{itemize}

\subsection{Time series basics}\label{time-series-basics}

\begin{Shaded}
\begin{Highlighting}[]
\FunctionTok{library}\NormalTok{(forecast)}
\NormalTok{ts\_obj }\OtherTok{\textless{}{-}} \FunctionTok{ts}\NormalTok{(x, }\AttributeTok{frequency =} \DecValTok{12}\NormalTok{)}
\FunctionTok{decompose}\NormalTok{(ts\_obj) }\SpecialCharTok{|\textgreater{}} \FunctionTok{plot}\NormalTok{()}
\FunctionTok{auto.arima}\NormalTok{(ts\_obj)}

\FunctionTok{library}\NormalTok{(changepoint)}
\FunctionTok{cpt.mean}\NormalTok{(x, }\AttributeTok{method =} \StringTok{"PELT"}\NormalTok{) }\SpecialCharTok{|\textgreater{}} \FunctionTok{plot}\NormalTok{()}
\end{Highlighting}
\end{Shaded}

\subsection{Decision rule}\label{decision-rule-9}

\begin{itemize}
\tightlist
\item
  Missing data \textgreater{} 5\% → at least sensitivity analysis;
  prefer MI.
\item
  Repeated measurements → mixed model, not repeated-measures ANOVA.
\item
  Binary outcome with clusters → GLMM (if subject effects matter) or GEE
  (if marginal effect is the estimand).
\item
  Time series with trend / seasonality → decompose before modelling.
\end{itemize}

\subsection{Common pitfalls}\label{common-pitfalls-9}

\begin{itemize}
\tightlist
\item
  Last-observation-carried-forward; usually biased.
\item
  Imputing after transformation rather than in the right scale.
\item
  Ignoring random effects when clusters are small in number (degrees of
  freedom correction).
\item
  Quoting GEE coefficients as if they were subject-specific.
\end{itemize}

\subsection{Further reading}\label{further-reading-9}

\begin{itemize}
\tightlist
\item
  van Buuren, \emph{Flexible Imputation of Missing Data}.
\item
  Fitzmaurice, Laird, Ware, \emph{Applied Longitudinal Analysis}.
\end{itemize}

\subsubsection{Causal inference methods\#\# Time-varying covariates /
landmarking}\label{causal-inference-methods-time-varying-covariates-landmarking}

\begin{itemize}
\tightlist
\item
  Counting-process data: split follow-up into intervals.
\item
  Landmark analysis: define a landmark time \(t^*\), condition on
  surviving to \(t^*\), classify by status at \(t^*\).
\end{itemize}

\begin{Shaded}
\begin{Highlighting}[]
\FunctionTok{library}\NormalTok{(survival)}
\FunctionTok{coxph}\NormalTok{(}\FunctionTok{Surv}\NormalTok{(tstart, tstop, event) }\SpecialCharTok{\textasciitilde{}}\NormalTok{ x, }\AttributeTok{data =}\NormalTok{ long)}
\end{Highlighting}
\end{Shaded}

\begin{itemize}
\tightlist
\item
  Immortal-time bias: exposed group can only be exposed if they live
  long enough to be exposed → spurious survival benefit.
\end{itemize}

\subsection{Competing risks}\label{competing-risks}

\begin{Shaded}
\begin{Highlighting}[]
\FunctionTok{library}\NormalTok{(tidycmprsk); }\FunctionTok{library}\NormalTok{(ggsurvfit)}
\NormalTok{cif }\OtherTok{\textless{}{-}} \FunctionTok{cuminc}\NormalTok{(}\FunctionTok{Surv}\NormalTok{(time, event\_f) }\SpecialCharTok{\textasciitilde{}}\NormalTok{ x, }\AttributeTok{data =}\NormalTok{ df)}
\NormalTok{ggsurvfit}\SpecialCharTok{::}\FunctionTok{ggcuminc}\NormalTok{(cif, }\AttributeTok{outcome =} \StringTok{"event1"}\NormalTok{)}

\CommentTok{\# Fine{-}Gray subdistribution hazard for event1}
\FunctionTok{crr}\NormalTok{(}\FunctionTok{Surv}\NormalTok{(time, event\_f) }\SpecialCharTok{\textasciitilde{}}\NormalTok{ x, }\AttributeTok{data =}\NormalTok{ df, }\AttributeTok{failcode =} \StringTok{"event1"}\NormalTok{)}
\end{Highlighting}
\end{Shaded}

\begin{itemize}
\tightlist
\item
  Cause-specific HR answers ``given you are event-free, what is the
  hazard for event \emph{k}?''
\item
  Subdistribution HR answers ``what is the hazard for event \emph{k} in
  the whole cohort, treating competing events as obstacles?''
\end{itemize}

\subsection{DAGs}\label{dags}

\begin{Shaded}
\begin{Highlighting}[]
\FunctionTok{library}\NormalTok{(dagitty); }\FunctionTok{library}\NormalTok{(ggdag)}
\NormalTok{dag }\OtherTok{\textless{}{-}} \FunctionTok{dagitty}\NormalTok{(}\StringTok{"dag \{ X {-}\textgreater{} Y ; C {-}\textgreater{} X ; C {-}\textgreater{} Y \}"}\NormalTok{)}
\FunctionTok{adjustmentSets}\NormalTok{(dag, }\AttributeTok{exposure =} \StringTok{"X"}\NormalTok{, }\AttributeTok{outcome =} \StringTok{"Y"}\NormalTok{)}
\FunctionTok{impliedConditionalIndependencies}\NormalTok{(dag)}
\end{Highlighting}
\end{Shaded}

\begin{itemize}
\tightlist
\item
  Adjust for confounders, \textbf{not} colliders.
\item
  Adjusting for a mediator estimates the direct effect only.
\end{itemize}

\subsection{Propensity scores / IPTW}\label{propensity-scores-iptw}

\begin{Shaded}
\begin{Highlighting}[]
\FunctionTok{library}\NormalTok{(MatchIt); }\FunctionTok{library}\NormalTok{(cobalt)}
\NormalTok{m }\OtherTok{\textless{}{-}} \FunctionTok{matchit}\NormalTok{(treat }\SpecialCharTok{\textasciitilde{}}\NormalTok{ x1 }\SpecialCharTok{+}\NormalTok{ x2, }\AttributeTok{data =}\NormalTok{ df, }\AttributeTok{method =} \StringTok{"nearest"}\NormalTok{)}
\FunctionTok{love.plot}\NormalTok{(m)   }\CommentTok{\# check balance: SMD \textless{} 0.1 is the goal}
\NormalTok{fit }\OtherTok{\textless{}{-}} \FunctionTok{lm}\NormalTok{(y }\SpecialCharTok{\textasciitilde{}}\NormalTok{ treat, }\AttributeTok{data =} \FunctionTok{match.data}\NormalTok{(m))}
\end{Highlighting}
\end{Shaded}

IPTW: weight by \(1/\hat{e}_i\) for treated, \(1/(1-\hat{e}_i)\) for
untreated. Stabilised weights stay closer to 1; trim extremes to avoid
outliers.

\subsection{G-methods, IV, DiD, RDD}\label{g-methods-iv-did-rdd}

{\def\LTcaptype{none} % do not increment counter
\begin{longtable}[]{@{}ll@{}}
\toprule\noalign{}
Method & Identifying assumption \\
\midrule\noalign{}
\endhead
\bottomrule\noalign{}
\endlastfoot
G-formula / IPTW & positivity + no unmeasured confounding \\
Instrumental variables & exclusion restriction + relevance \\
Difference-in-differences & parallel trends \\
Regression discontinuity & continuity at the cutoff \\
\end{longtable}
}

\subsection{Heterogeneous treatment
effects}\label{heterogeneous-treatment-effects}

\begin{itemize}
\tightlist
\item
  Causal forest (\texttt{grf::causal\_forest}) or meta-learners (S-, T-,
  X-).
\item
  Report conditional ATEs with adjusted intervals; avoid spurious
  subgroup p-values.
\end{itemize}

\subsection{Decision rule}\label{decision-rule-10}

\begin{itemize}
\tightlist
\item
  Treatment assigned over time → time-varying Cox.
\item
  Multiple competing events → cumulative incidence, Fine-Gray.
\item
  Observational causal question → DAG first, choose identification
  strategy second.
\item
  Want the whole distribution of treatment effect? → causal forest.
\end{itemize}

\subsection{Common pitfalls}\label{common-pitfalls-10}

\begin{itemize}
\tightlist
\item
  Adjusting for a post-treatment variable (induces collider bias).
\item
  Quoting HR from a Cox model that violates PH.
\item
  Matching but never checking balance.
\item
  Calling IV estimates ``generalisable'' when they are local to
  compliers.
\end{itemize}

\subsection{Further reading}\label{further-reading-10}

\begin{itemize}
\tightlist
\item
  Hernán \& Robins, \emph{Causal Inference: What If}.
\item
  Therneau \& Grambsch, \emph{Modeling Survival Data}.
\end{itemize}

\subsubsection{Synthesis and pre-registration\#\# Systematic review
workflow}\label{synthesis-and-pre-registration-systematic-review-workflow}

\begin{enumerate}
\def\labelenumi{\arabic{enumi}.}
\tightlist
\item
  PICO: Population, Intervention, Comparator, Outcome.
\item
  Register protocol on PROSPERO.
\item
  Build search across ≥ 2 databases; peer-review with a librarian.
\item
  Two-screener title/abstract and full-text.
\item
  Extract, appraise (RoB 2 / ROBINS-I), synthesise.
\item
  Report with the PRISMA 2020 checklist and flow diagram.
\end{enumerate}

\subsection{Meta-analysis}\label{meta-analysis}

\begin{Shaded}
\begin{Highlighting}[]
\FunctionTok{library}\NormalTok{(metafor)}
\NormalTok{dat }\OtherTok{\textless{}{-}} \FunctionTok{escalc}\NormalTok{(}\AttributeTok{measure =} \StringTok{"RR"}\NormalTok{,}
              \AttributeTok{ai =}\NormalTok{ tpos, }\AttributeTok{bi =}\NormalTok{ tneg, }\AttributeTok{ci =}\NormalTok{ cpos, }\AttributeTok{di =}\NormalTok{ cneg,}
              \AttributeTok{data =}\NormalTok{ studies)}
\NormalTok{fit }\OtherTok{\textless{}{-}} \FunctionTok{rma}\NormalTok{(yi, vi, }\AttributeTok{data =}\NormalTok{ dat, }\AttributeTok{method =} \StringTok{"REML"}\NormalTok{)}
\NormalTok{fit}
\FunctionTok{forest}\NormalTok{(fit, }\AttributeTok{header =} \ConstantTok{TRUE}\NormalTok{)}
\FunctionTok{funnel}\NormalTok{(fit); }\FunctionTok{regtest}\NormalTok{(fit)}
\end{Highlighting}
\end{Shaded}

{\def\LTcaptype{none} % do not increment counter
\begin{longtable}[]{@{}ll@{}}
\toprule\noalign{}
Statistic & Meaning \\
\midrule\noalign{}
\endhead
\bottomrule\noalign{}
\endlastfoot
\(\tau^2\) & between-study variance \\
\(I^2\) & \% of total variance due to heterogeneity \\
Egger's test & funnel-plot asymmetry (small-study effects) \\
\end{longtable}
}

Fixed-effect only when heterogeneity is effectively zero. Default to
random effects.

\subsection{Network meta-analysis}\label{network-meta-analysis}

\begin{Shaded}
\begin{Highlighting}[]
\FunctionTok{library}\NormalTok{(netmeta)}
\NormalTok{net }\OtherTok{\textless{}{-}} \FunctionTok{netmeta}\NormalTok{(TE, seTE, treat1, treat2, studlab,}
               \AttributeTok{data =}\NormalTok{ dat, }\AttributeTok{sm =} \StringTok{"MD"}\NormalTok{, }\AttributeTok{random =} \ConstantTok{TRUE}\NormalTok{)}
\FunctionTok{netgraph}\NormalTok{(net)}
\FunctionTok{netrank}\NormalTok{(net, }\AttributeTok{small.values =} \StringTok{"good"}\NormalTok{)    }\CommentTok{\# SUCRA ranking}
\FunctionTok{decomp.design}\NormalTok{(net)                     }\CommentTok{\# inconsistency}
\end{Highlighting}
\end{Shaded}

Transitivity + consistency are the core assumptions.

\subsection{Infectious disease --- SIR /
SEIR}\label{infectious-disease-sir-seir}

\begin{Shaded}
\begin{Highlighting}[]
\FunctionTok{library}\NormalTok{(deSolve)}
\NormalTok{sir }\OtherTok{\textless{}{-}} \ControlFlowTok{function}\NormalTok{(t, y, p) }\FunctionTok{with}\NormalTok{(}\FunctionTok{as.list}\NormalTok{(}\FunctionTok{c}\NormalTok{(y, p)), \{}
\NormalTok{  N }\OtherTok{\textless{}{-}}\NormalTok{ S }\SpecialCharTok{+}\NormalTok{ I }\SpecialCharTok{+}\NormalTok{ R}
  \FunctionTok{list}\NormalTok{(}\FunctionTok{c}\NormalTok{(}\SpecialCharTok{{-}}\NormalTok{beta }\SpecialCharTok{*}\NormalTok{ S }\SpecialCharTok{*}\NormalTok{ I }\SpecialCharTok{/}\NormalTok{ N,}
\NormalTok{          beta }\SpecialCharTok{*}\NormalTok{ S }\SpecialCharTok{*}\NormalTok{ I }\SpecialCharTok{/}\NormalTok{ N }\SpecialCharTok{{-}}\NormalTok{ gamma }\SpecialCharTok{*}\NormalTok{ I,}
\NormalTok{          gamma }\SpecialCharTok{*}\NormalTok{ I))}
\NormalTok{\})}
\NormalTok{out }\OtherTok{\textless{}{-}} \FunctionTok{ode}\NormalTok{(}\FunctionTok{c}\NormalTok{(}\AttributeTok{S =} \DecValTok{999}\NormalTok{, }\AttributeTok{I =} \DecValTok{1}\NormalTok{, }\AttributeTok{R =} \DecValTok{0}\NormalTok{),}
           \AttributeTok{times =} \DecValTok{0}\SpecialCharTok{:}\DecValTok{120}\NormalTok{,}
           \AttributeTok{func =}\NormalTok{ sir,}
           \AttributeTok{parms =} \FunctionTok{c}\NormalTok{(}\AttributeTok{beta =} \FloatTok{0.3}\NormalTok{, }\AttributeTok{gamma =} \FloatTok{0.1}\NormalTok{))}
\end{Highlighting}
\end{Shaded}

\(R_0 = \beta / \gamma\). Herd-immunity threshold \(\approx 1 - 1/R_0\).

\subsection{Pre-registration and SAP}\label{pre-registration-and-sap}

A pre-registration specifies, before data analysis:

\begin{itemize}
\tightlist
\item
  Primary outcome + exact estimand.
\item
  Analysis model (fixed terms, adjustments, interactions).
\item
  Handling of missing data.
\item
  Multiplicity control.
\item
  Sample-size justification.
\end{itemize}

OSF and ClinicalTrials.gov are the two common registries.

\subsection{Decision rule}\label{decision-rule-11}

\begin{itemize}
\tightlist
\item
  More than a handful of studies on a question → plan a systematic
  review.
\item
  Multiple comparators of interest → network meta-analysis.
\item
  Outbreak modelling → start with SIR, extend only as data justify.
\item
  Confirmatory analysis → pre-registered; exploratory analyses are
  labelled as such in the manuscript.
\end{itemize}

\subsection{Common pitfalls}\label{common-pitfalls-11}

\begin{itemize}
\tightlist
\item
  Declaring a review ``systematic'' without a pre-registered protocol.
\item
  Pooling studies with fundamentally different populations / outcomes.
\item
  Interpreting SUCRA rankings as if they were certain.
\item
  Calling a compartmental model a ``forecast'' without uncertainty.
\end{itemize}

\subsection{Further reading}\label{further-reading-11}

\begin{itemize}
\tightlist
\item
  Higgins et al., \emph{Cochrane Handbook for Systematic Reviews}.
\item
  Keeling \& Rohani, \emph{Modeling Infectious Diseases in Humans and
  Animals}.
\item
  Nosek et al.~(2018), \emph{The preregistration revolution}.
\end{itemize}

\subsection{Machine Learning \& High-Dimensional
Data}\label{machine-learning-high-dimensional-data}

\subsubsection{Cross-validation, regularisation, and unsupervised\#\#
Cross-validation}\label{cross-validation-regularisation-and-unsupervised-cross-validation}

{\def\LTcaptype{none} % do not increment counter
\begin{longtable}[]{@{}ll@{}}
\toprule\noalign{}
Scheme & Use \\
\midrule\noalign{}
\endhead
\bottomrule\noalign{}
\endlastfoot
k-fold (k = 5 or 10) & baseline, fast \\
Repeated k-fold & reduces Monte Carlo noise \\
Leave-one-out & small n, high variance \\
Stratified k-fold & imbalanced class labels \\
Grouped / blocked & correlated units (patients, sites) \\
Nested & honest evaluation of tuned models \\
\end{longtable}
}

Nested CV is the only defensible way to estimate generalisation error
when you are also tuning hyperparameters.

\subsection{Regularisation}\label{regularisation}

\begin{Shaded}
\begin{Highlighting}[]
\FunctionTok{library}\NormalTok{(glmnet)}
\NormalTok{x }\OtherTok{\textless{}{-}} \FunctionTok{model.matrix}\NormalTok{(y }\SpecialCharTok{\textasciitilde{}}\NormalTok{ . }\SpecialCharTok{{-}} \DecValTok{1}\NormalTok{, df); y }\OtherTok{\textless{}{-}}\NormalTok{ df}\SpecialCharTok{$}\NormalTok{y}
\NormalTok{fit }\OtherTok{\textless{}{-}} \FunctionTok{cv.glmnet}\NormalTok{(x, y, }\AttributeTok{alpha =} \DecValTok{1}\NormalTok{)     }\CommentTok{\# alpha = 1 = lasso; 0 = ridge}
\FunctionTok{coef}\NormalTok{(fit, }\AttributeTok{s =} \StringTok{"lambda.1se"}\NormalTok{)           }\CommentTok{\# sparse solution}
\FunctionTok{plot}\NormalTok{(fit)                              }\CommentTok{\# CV curve}
\end{Highlighting}
\end{Shaded}

\begin{itemize}
\tightlist
\item
  \texttt{lambda.min}: empirical minimum of CV error.
\item
  \texttt{lambda.1se}: more regularised, simpler, within 1 SE of min.
\item
  Elastic net: \texttt{alpha} between 0 and 1 blends ridge and lasso.
\end{itemize}

\subsection{Multivariate workhorses}\label{multivariate-workhorses}

{\def\LTcaptype{none} % do not increment counter
\begin{longtable}[]{@{}ll@{}}
\toprule\noalign{}
Method & Purpose \\
\midrule\noalign{}
\endhead
\bottomrule\noalign{}
\endlastfoot
PCA & variance → low-d summary; \texttt{prcomp(x,\ scale.\ =\ TRUE)} \\
FA & latent-variable modelling; \texttt{psych::fa} \\
CCA & correlation structure between two sets \\
LDA & supervised low-d projection; \texttt{MASS::lda} \\
\end{longtable}
}

\begin{Shaded}
\begin{Highlighting}[]
\NormalTok{pc }\OtherTok{\textless{}{-}} \FunctionTok{prcomp}\NormalTok{(x, }\AttributeTok{scale. =} \ConstantTok{TRUE}\NormalTok{)}
\FunctionTok{summary}\NormalTok{(pc)   }\CommentTok{\# proportion of variance per PC}
\end{Highlighting}
\end{Shaded}

Scale before PCA when variables are on different units.

\subsection{Clustering}\label{clustering}

\begin{Shaded}
\begin{Highlighting}[]
\NormalTok{km }\OtherTok{\textless{}{-}} \FunctionTok{kmeans}\NormalTok{(}\FunctionTok{scale}\NormalTok{(x), }\AttributeTok{centers =} \DecValTok{4}\NormalTok{, }\AttributeTok{nstart =} \DecValTok{25}\NormalTok{)}
\NormalTok{hc }\OtherTok{\textless{}{-}} \FunctionTok{hclust}\NormalTok{(}\FunctionTok{dist}\NormalTok{(}\FunctionTok{scale}\NormalTok{(x)), }\AttributeTok{method =} \StringTok{"ward.D2"}\NormalTok{)}
\FunctionTok{library}\NormalTok{(mclust); bic }\OtherTok{\textless{}{-}} \FunctionTok{mclustBIC}\NormalTok{(x)      }\CommentTok{\# model{-}based, BIC{-}selected}
\end{Highlighting}
\end{Shaded}

\begin{itemize}
\tightlist
\item
  Choose k by silhouette, gap statistic, or BIC.
\item
  Scale first. K-means assumes round clusters of similar size.
\end{itemize}

\subsection{Non-linear embeddings}\label{non-linear-embeddings}

\begin{Shaded}
\begin{Highlighting}[]
\FunctionTok{library}\NormalTok{(uwot)}
\NormalTok{umap\_xy }\OtherTok{\textless{}{-}} \FunctionTok{umap}\NormalTok{(}\FunctionTok{scale}\NormalTok{(x), }\AttributeTok{n\_neighbors =} \DecValTok{15}\NormalTok{, }\AttributeTok{min\_dist =} \FloatTok{0.1}\NormalTok{)}

\FunctionTok{library}\NormalTok{(Rtsne)}
\NormalTok{ts }\OtherTok{\textless{}{-}} \FunctionTok{Rtsne}\NormalTok{(}\FunctionTok{scale}\NormalTok{(x), }\AttributeTok{perplexity =} \DecValTok{30}\NormalTok{)}\SpecialCharTok{$}\NormalTok{Y}
\end{Highlighting}
\end{Shaded}

Embeddings are \emph{visualisations}, not distances you should input
into downstream stats. Global geometry is often distorted.

\subsection{Decision rule}\label{decision-rule-12}

\begin{itemize}
\tightlist
\item
  Lots of predictors, small n → lasso; report \texttt{lambda.1se}.
\item
  Exploring structure → PCA + UMAP on scaled data.
\item
  Clustering for biology → model-based + sensitivity to k.
\item
  Tuning anything → nested CV.
\end{itemize}

\subsection{Common pitfalls}\label{common-pitfalls-12}

\begin{itemize}
\tightlist
\item
  Selecting features on full data, then estimating error by k-fold.
\item
  PCA before scaling when features are on different units.
\item
  Interpreting UMAP distances as real distances.
\item
  Using AUC on rank-ordered hold-out error as if it were CV error.
\end{itemize}

\subsection{Further reading}\label{further-reading-12}

\begin{itemize}
\tightlist
\item
  Hastie, Tibshirani, Friedman, \emph{The Elements of Statistical
  Learning}.
\item
  James et al., \emph{ISLR2}.
\end{itemize}

\subsubsection{Trees, boosting, and interpretability\#\# Tree-based
models}\label{trees-boosting-and-interpretability-tree-based-models}

{\def\LTcaptype{none} % do not increment counter
\begin{longtable}[]{@{}
  >{\raggedright\arraybackslash}p{(\linewidth - 4\tabcolsep) * \real{0.3333}}
  >{\raggedright\arraybackslash}p{(\linewidth - 4\tabcolsep) * \real{0.3333}}
  >{\raggedright\arraybackslash}p{(\linewidth - 4\tabcolsep) * \real{0.3333}}@{}}
\toprule\noalign{}
\begin{minipage}[b]{\linewidth}\raggedright
Model
\end{minipage} & \begin{minipage}[b]{\linewidth}\raggedright
R
\end{minipage} & \begin{minipage}[b]{\linewidth}\raggedright
Strengths
\end{minipage} \\
\midrule\noalign{}
\endhead
\bottomrule\noalign{}
\endlastfoot
CART & \texttt{rpart::rpart} & interpretable single tree \\
Random forest & \texttt{ranger::ranger} & robust, OOB error, variable
importance \\
XGBoost & \texttt{xgboost::xgb.train} & top tabular performance, careful
tuning \\
LightGBM & \texttt{lightgbm::lgb.train} & fast on large data \\
\end{longtable}
}

\begin{Shaded}
\begin{Highlighting}[]
\FunctionTok{library}\NormalTok{(ranger)}
\NormalTok{fit }\OtherTok{\textless{}{-}} \FunctionTok{ranger}\NormalTok{(y }\SpecialCharTok{\textasciitilde{}}\NormalTok{ ., }\AttributeTok{data =}\NormalTok{ df, }\AttributeTok{importance =} \StringTok{"permutation"}\NormalTok{,}
              \AttributeTok{num.trees =} \DecValTok{500}\NormalTok{, }\AttributeTok{mtry =} \FunctionTok{floor}\NormalTok{(}\FunctionTok{sqrt}\NormalTok{(}\FunctionTok{ncol}\NormalTok{(df) }\SpecialCharTok{{-}} \DecValTok{1}\NormalTok{)))}
\NormalTok{fit}\SpecialCharTok{$}\NormalTok{prediction.error   }\CommentTok{\# OOB error}
\end{Highlighting}
\end{Shaded}

\subsection{Interpretability}\label{interpretability}

\begin{Shaded}
\begin{Highlighting}[]
\FunctionTok{library}\NormalTok{(DALEX); }\FunctionTok{library}\NormalTok{(iml)}

\NormalTok{ex  }\OtherTok{\textless{}{-}} \FunctionTok{explain}\NormalTok{(fit, }\AttributeTok{data =}\NormalTok{ df }\SpecialCharTok{|\textgreater{}} \FunctionTok{select}\NormalTok{(}\SpecialCharTok{{-}}\NormalTok{y), }\AttributeTok{y =}\NormalTok{ df}\SpecialCharTok{$}\NormalTok{y)}
\NormalTok{ip  }\OtherTok{\textless{}{-}} \FunctionTok{model\_parts}\NormalTok{(ex)           }\CommentTok{\# permutation importance}
\NormalTok{pd  }\OtherTok{\textless{}{-}} \FunctionTok{model\_profile}\NormalTok{(ex, }\AttributeTok{variables =} \StringTok{"x1"}\NormalTok{)   }\CommentTok{\# partial dependence}
\NormalTok{sh  }\OtherTok{\textless{}{-}}\NormalTok{ Shapley}\SpecialCharTok{$}\FunctionTok{new}\NormalTok{(Predictor}\SpecialCharTok{$}\FunctionTok{new}\NormalTok{(fit, }\AttributeTok{data =}\NormalTok{ df, }\AttributeTok{y =}\NormalTok{ df}\SpecialCharTok{$}\NormalTok{y),}
                   \AttributeTok{x.interest =}\NormalTok{ df[}\DecValTok{1}\NormalTok{, ])     }\CommentTok{\# SHAP}
\end{Highlighting}
\end{Shaded}

PDP assumes feature independence; if features are correlated, use ALE
plots.

\subsection{\texorpdfstring{Tabular neural networks with
\texttt{torch}}{Tabular neural networks with torch}}\label{tabular-neural-networks-with-torch}

\begin{Shaded}
\begin{Highlighting}[]
\CommentTok{\# sketch only — typically in a loop on GPU}
\FunctionTok{library}\NormalTok{(torch)}
\NormalTok{nn }\OtherTok{\textless{}{-}} \FunctionTok{nn\_sequential}\NormalTok{(}
  \FunctionTok{nn\_linear}\NormalTok{(p, }\DecValTok{64}\NormalTok{), }\FunctionTok{nn\_relu}\NormalTok{(),}
  \FunctionTok{nn\_linear}\NormalTok{(}\DecValTok{64}\NormalTok{, }\DecValTok{1}\NormalTok{)}
\NormalTok{)}
\NormalTok{opt }\OtherTok{\textless{}{-}} \FunctionTok{optim\_adam}\NormalTok{(nn}\SpecialCharTok{$}\NormalTok{parameters)}
\end{Highlighting}
\end{Shaded}

On small biomedical tabular data, boosted trees usually beat NN. Use NN
when you need end-to-end training with images, sequences, or text.

\subsection{tidymodels pipeline}\label{tidymodels-pipeline}

\begin{Shaded}
\begin{Highlighting}[]
\FunctionTok{library}\NormalTok{(tidymodels)}
\NormalTok{rec  }\OtherTok{\textless{}{-}} \FunctionTok{recipe}\NormalTok{(y }\SpecialCharTok{\textasciitilde{}}\NormalTok{ ., }\AttributeTok{data =}\NormalTok{ df) }\SpecialCharTok{|\textgreater{}}
  \FunctionTok{step\_normalize}\NormalTok{(}\FunctionTok{all\_numeric\_predictors}\NormalTok{()) }\SpecialCharTok{|\textgreater{}}
  \FunctionTok{step\_dummy}\NormalTok{(}\FunctionTok{all\_nominal\_predictors}\NormalTok{())}

\NormalTok{mod  }\OtherTok{\textless{}{-}} \FunctionTok{rand\_forest}\NormalTok{(}\AttributeTok{mtry =} \FunctionTok{tune}\NormalTok{(), }\AttributeTok{trees =} \DecValTok{500}\NormalTok{) }\SpecialCharTok{|\textgreater{}}
  \FunctionTok{set\_engine}\NormalTok{(}\StringTok{"ranger"}\NormalTok{) }\SpecialCharTok{|\textgreater{}} \FunctionTok{set\_mode}\NormalTok{(}\StringTok{"classification"}\NormalTok{)}

\NormalTok{wf   }\OtherTok{\textless{}{-}} \FunctionTok{workflow}\NormalTok{() }\SpecialCharTok{|\textgreater{}} \FunctionTok{add\_recipe}\NormalTok{(rec) }\SpecialCharTok{|\textgreater{}} \FunctionTok{add\_model}\NormalTok{(mod)}
\NormalTok{cv   }\OtherTok{\textless{}{-}} \FunctionTok{vfold\_cv}\NormalTok{(df, }\AttributeTok{v =} \DecValTok{5}\NormalTok{, }\AttributeTok{strata =}\NormalTok{ y)}
\NormalTok{tuned }\OtherTok{\textless{}{-}} \FunctionTok{tune\_grid}\NormalTok{(wf, cv, }\AttributeTok{grid =} \DecValTok{20}\NormalTok{)}
\FunctionTok{collect\_metrics}\NormalTok{(tuned)}
\end{Highlighting}
\end{Shaded}

\subsection{Decision rule}\label{decision-rule-13}

\begin{itemize}
\tightlist
\item
  Tabular, \textless{} 100k rows → gradient-boosted trees.
\item
  Need explanations → SHAP + PDP, validate with held-out.
\item
  Images / sequences → CNN / transformer, use \texttt{torch}.
\item
  Reproducibility → \texttt{tidymodels} pipeline + fixed seed + locked
  recipe.
\end{itemize}

\subsection{Common pitfalls}\label{common-pitfalls-13}

\begin{itemize}
\tightlist
\item
  Reporting feature importance from a model trained on the full data.
\item
  Interpreting PDPs when features are strongly correlated.
\item
  Claiming NN superiority without CV on the same split as a tree model.
\item
  Forgetting to set a seed for any tune / split operation.
\end{itemize}

\subsection{Further reading}\label{further-reading-13}

\begin{itemize}
\tightlist
\item
  Biecek \& Burzykowski, \emph{Explanatory Model Analysis}.
\item
  Molnar, \emph{Interpretable Machine Learning}.
\end{itemize}

\subsubsection{Bayesian and survival ML\#\# Bayesian
thinking}\label{bayesian-and-survival-ml-bayesian-thinking}

\(\underbrace{p(\theta \mid y)}_{\text{posterior}} \propto
 \underbrace{p(y \mid \theta)}_{\text{likelihood}}
 \underbrace{p(\theta)}_{\text{prior}}\)

\begin{itemize}
\tightlist
\item
  Priors are part of the model --- specify and defend them.
\item
  Posterior summaries: mean, median, 95\% credible interval.
\item
  No \emph{p}-values; use posterior probability of direction, Bayes
  factors, or LOO.
\end{itemize}

\subsection{\texorpdfstring{\texttt{brms} /
Stan}{brms / Stan}}\label{brms-stan}

\begin{Shaded}
\begin{Highlighting}[]
\FunctionTok{library}\NormalTok{(brms)}
\NormalTok{fit }\OtherTok{\textless{}{-}} \FunctionTok{brm}\NormalTok{(y }\SpecialCharTok{\textasciitilde{}}\NormalTok{ x }\SpecialCharTok{+}\NormalTok{ (}\DecValTok{1} \SpecialCharTok{|}\NormalTok{ group), }\AttributeTok{data =}\NormalTok{ df,}
           \AttributeTok{prior =} \FunctionTok{c}\NormalTok{(}\FunctionTok{prior}\NormalTok{(}\FunctionTok{normal}\NormalTok{(}\DecValTok{0}\NormalTok{, }\DecValTok{1}\NormalTok{), }\AttributeTok{class =} \StringTok{"b"}\NormalTok{),}
                     \FunctionTok{prior}\NormalTok{(}\FunctionTok{exponential}\NormalTok{(}\DecValTok{1}\NormalTok{), }\AttributeTok{class =} \StringTok{"sd"}\NormalTok{)),}
           \AttributeTok{chains =} \DecValTok{4}\NormalTok{, }\AttributeTok{iter =} \DecValTok{2000}\NormalTok{, }\AttributeTok{seed =} \DecValTok{42}\NormalTok{)}

\FunctionTok{summary}\NormalTok{(fit)}
\FunctionTok{loo}\NormalTok{(fit)               }\CommentTok{\# leave{-}one{-}out cross{-}validation}
\FunctionTok{pp\_check}\NormalTok{(fit)           }\CommentTok{\# posterior predictive checks}
\end{Highlighting}
\end{Shaded}

Prior predictive check before fitting: simulate \(y\) from the prior and
confirm the implications are plausible.

\subsection{Biomarker statistics}\label{biomarker-statistics}

{\def\LTcaptype{none} % do not increment counter
\begin{longtable}[]{@{}
  >{\raggedright\arraybackslash}p{(\linewidth - 2\tabcolsep) * \real{0.5000}}
  >{\raggedright\arraybackslash}p{(\linewidth - 2\tabcolsep) * \real{0.5000}}@{}}
\toprule\noalign{}
\begin{minipage}[b]{\linewidth}\raggedright
Question
\end{minipage} & \begin{minipage}[b]{\linewidth}\raggedright
Statistic
\end{minipage} \\
\midrule\noalign{}
\endhead
\bottomrule\noalign{}
\endlastfoot
Does a biomarker classify? & AUC, cut-point (Youden's index) \\
Does it add over an existing model? & ΔAUC, NRI, IDI, decision curves \\
Does it move prognosis? & Calibration plus discrimination \\
Is a cut-off reproducible? & Bootstrap CI on the cut-off \\
\end{longtable}
}

\begin{Shaded}
\begin{Highlighting}[]
\NormalTok{pROC}\SpecialCharTok{::}\FunctionTok{roc}\NormalTok{(y, biomarker) }\SpecialCharTok{|\textgreater{}}
\NormalTok{  pROC}\SpecialCharTok{::}\FunctionTok{coords}\NormalTok{(}\StringTok{"best"}\NormalTok{, }\AttributeTok{best.method =} \StringTok{"youden"}\NormalTok{)}
\end{Highlighting}
\end{Shaded}

\subsection{Survival ML}\label{survival-ml}

{\def\LTcaptype{none} % do not increment counter
\begin{longtable}[]{@{}
  >{\raggedright\arraybackslash}p{(\linewidth - 2\tabcolsep) * \real{0.5000}}
  >{\raggedright\arraybackslash}p{(\linewidth - 2\tabcolsep) * \real{0.5000}}@{}}
\toprule\noalign{}
\begin{minipage}[b]{\linewidth}\raggedright
Model
\end{minipage} & \begin{minipage}[b]{\linewidth}\raggedright
R
\end{minipage} \\
\midrule\noalign{}
\endhead
\bottomrule\noalign{}
\endlastfoot
Random survival forest &
\texttt{randomForestSRC::rfsrc(Surv(...)\ \textasciitilde{}\ .,\ data)} \\
Gradient-boosted Cox & \texttt{xgboost} with
\texttt{objective\ =\ "survival:cox"} \\
DeepSurv (conceptual) & \texttt{torch} with partial-likelihood loss \\
\end{longtable}
}

Evaluate with \textbf{time-dependent AUC}, \textbf{integrated Brier
score}, and \textbf{IPA} (index of prediction accuracy = 1 −
Brier(model) / Brier(null)).

\begin{Shaded}
\begin{Highlighting}[]
\FunctionTok{library}\NormalTok{(timeROC)}
\NormalTok{r }\OtherTok{\textless{}{-}} \FunctionTok{timeROC}\NormalTok{(}\AttributeTok{T =}\NormalTok{ df}\SpecialCharTok{$}\NormalTok{time, }\AttributeTok{delta =}\NormalTok{ df}\SpecialCharTok{$}\NormalTok{event,}
             \AttributeTok{marker =}\NormalTok{ pred, }\AttributeTok{cause =} \DecValTok{1}\NormalTok{, }\AttributeTok{times =} \FunctionTok{c}\NormalTok{(}\DecValTok{365}\NormalTok{, }\DecValTok{730}\NormalTok{))}
\end{Highlighting}
\end{Shaded}

\subsection{External validation}\label{external-validation}

\begin{itemize}
\tightlist
\item
  Never trust the apparent performance on training data.
\item
  Minimum: split-sample or CV on internal data.
\item
  Better: external validation in a second cohort.
\item
  Report calibration slope / intercept, discrimination, and decision
  curves.
\end{itemize}

\subsection{Decision rule}\label{decision-rule-14}

\begin{itemize}
\tightlist
\item
  Rare outcome with need for uncertainty → Bayesian, not bootstrap of
  MLE.
\item
  New biomarker → NRI + decision curve, not just ΔAUC.
\item
  Survival prediction → time-dependent Brier / IPA, not global AUC.
\item
  Any clinical claim → external validation before publication.
\end{itemize}

\subsection{Common pitfalls}\label{common-pitfalls-14}

\begin{itemize}
\tightlist
\item
  Reporting a posterior with default flat priors on unscaled predictors.
\item
  Selecting the Youden cut-off on the full data and using the same data
  to evaluate sensitivity / specificity.
\item
  Quoting C-statistic at a single time point and calling it survival ML.
\item
  Publishing a prediction model without TRIPOD-compliant reporting.
\end{itemize}

\subsection{Further reading}\label{further-reading-14}

\begin{itemize}
\tightlist
\item
  Gelman et al., \emph{Bayesian Data Analysis}, 3e.
\item
  Royston \& Altman, \emph{External validation of a Cox prognostic
  model}.
\end{itemize}

\subsubsection{Omics, FDR, and TRIPOD-AI\#\# Bulk
RNA-seq}\label{omics-fdr-and-tripod-ai-bulk-rna-seq}

\begin{Shaded}
\begin{Highlighting}[]
\FunctionTok{library}\NormalTok{(DESeq2)}
\NormalTok{dds }\OtherTok{\textless{}{-}} \FunctionTok{DESeqDataSetFromMatrix}\NormalTok{(counts, coldata, }\SpecialCharTok{\textasciitilde{}}\NormalTok{ condition)}
\NormalTok{dds }\OtherTok{\textless{}{-}} \FunctionTok{DESeq}\NormalTok{(dds)}
\NormalTok{res }\OtherTok{\textless{}{-}} \FunctionTok{results}\NormalTok{(dds, }\AttributeTok{contrast =} \FunctionTok{c}\NormalTok{(}\StringTok{"condition"}\NormalTok{, }\StringTok{"B"}\NormalTok{, }\StringTok{"A"}\NormalTok{))}
\end{Highlighting}
\end{Shaded}

\begin{Shaded}
\begin{Highlighting}[]
\FunctionTok{library}\NormalTok{(edgeR)}
\NormalTok{dge }\OtherTok{\textless{}{-}} \FunctionTok{DGEList}\NormalTok{(counts, }\AttributeTok{group =}\NormalTok{ coldata}\SpecialCharTok{$}\NormalTok{condition)}
\NormalTok{dge }\OtherTok{\textless{}{-}} \FunctionTok{calcNormFactors}\NormalTok{(dge)}
\NormalTok{design }\OtherTok{\textless{}{-}} \FunctionTok{model.matrix}\NormalTok{(}\SpecialCharTok{\textasciitilde{}}\NormalTok{ condition, coldata)}
\NormalTok{fit }\OtherTok{\textless{}{-}} \FunctionTok{glmQLFit}\NormalTok{(}\FunctionTok{estimateDisp}\NormalTok{(dge, design), design)}
\NormalTok{qlf }\OtherTok{\textless{}{-}} \FunctionTok{glmQLFTest}\NormalTok{(fit, }\AttributeTok{coef =} \DecValTok{2}\NormalTok{)}
\FunctionTok{topTags}\NormalTok{(qlf)}
\end{Highlighting}
\end{Shaded}

{\def\LTcaptype{none} % do not increment counter
\begin{longtable}[]{@{}
  >{\raggedright\arraybackslash}p{(\linewidth - 4\tabcolsep) * \real{0.3333}}
  >{\raggedright\arraybackslash}p{(\linewidth - 4\tabcolsep) * \real{0.3333}}
  >{\raggedright\arraybackslash}p{(\linewidth - 4\tabcolsep) * \real{0.3333}}@{}}
\toprule\noalign{}
\begin{minipage}[b]{\linewidth}\raggedright
Package
\end{minipage} & \begin{minipage}[b]{\linewidth}\raggedright
Distribution
\end{minipage} & \begin{minipage}[b]{\linewidth}\raggedright
Notes
\end{minipage} \\
\midrule\noalign{}
\endhead
\bottomrule\noalign{}
\endlastfoot
DESeq2 & negative binomial & shrinks LFCs; standard for small cohorts \\
edgeR & negative binomial & QL-F test has good type-I control \\
limma-voom & mean-variance weights & fast on very large studies \\
\end{longtable}
}

\subsection{Enrichment}\label{enrichment}

\begin{Shaded}
\begin{Highlighting}[]
\FunctionTok{library}\NormalTok{(fgsea)}
\FunctionTok{fgsea}\NormalTok{(pathways, }\AttributeTok{stats =}\NormalTok{ stat\_vector, }\AttributeTok{nperm =} \DecValTok{10000}\NormalTok{) }\SpecialCharTok{|\textgreater{}}
  \FunctionTok{arrange}\NormalTok{(padj)}

\CommentTok{\# Over{-}representation with a hypergeometric test}
\FunctionTok{phyper}\NormalTok{(k }\SpecialCharTok{{-}} \DecValTok{1}\NormalTok{, K, N }\SpecialCharTok{{-}}\NormalTok{ K, n, }\AttributeTok{lower.tail =} \ConstantTok{FALSE}\NormalTok{)}
\end{Highlighting}
\end{Shaded}

Always pre-rank by a \emph{signed} statistic (log-FC × -log10 p), not by
raw p alone.

\subsection{Single-cell RNA-seq (Seurat
pattern)}\label{single-cell-rna-seq-seurat-pattern}

\begin{Shaded}
\begin{Highlighting}[]
\FunctionTok{library}\NormalTok{(Seurat)}
\NormalTok{so }\OtherTok{\textless{}{-}} \FunctionTok{CreateSeuratObject}\NormalTok{(counts) }\SpecialCharTok{|\textgreater{}}
  \FunctionTok{NormalizeData}\NormalTok{() }\SpecialCharTok{|\textgreater{}}
  \FunctionTok{FindVariableFeatures}\NormalTok{() }\SpecialCharTok{|\textgreater{}}
  \FunctionTok{ScaleData}\NormalTok{() }\SpecialCharTok{|\textgreater{}}
  \FunctionTok{RunPCA}\NormalTok{() }\SpecialCharTok{|\textgreater{}}
  \FunctionTok{FindNeighbors}\NormalTok{() }\SpecialCharTok{|\textgreater{}}
  \FunctionTok{FindClusters}\NormalTok{(}\AttributeTok{resolution =} \FloatTok{0.4}\NormalTok{) }\SpecialCharTok{|\textgreater{}}
  \FunctionTok{RunUMAP}\NormalTok{(}\AttributeTok{dims =} \DecValTok{1}\SpecialCharTok{:}\DecValTok{20}\NormalTok{)}
\FunctionTok{DimPlot}\NormalTok{(so)}
\end{Highlighting}
\end{Shaded}

\subsection{FDR, knockoffs,
replication}\label{fdr-knockoffs-replication}

{\def\LTcaptype{none} % do not increment counter
\begin{longtable}[]{@{}
  >{\raggedright\arraybackslash}p{(\linewidth - 2\tabcolsep) * \real{0.5000}}
  >{\raggedright\arraybackslash}p{(\linewidth - 2\tabcolsep) * \real{0.5000}}@{}}
\toprule\noalign{}
\begin{minipage}[b]{\linewidth}\raggedright
Correction
\end{minipage} & \begin{minipage}[b]{\linewidth}\raggedright
When
\end{minipage} \\
\midrule\noalign{}
\endhead
\bottomrule\noalign{}
\endlastfoot
Bonferroni & few planned tests, strict control \\
Benjamini-Hochberg & many tests, accept proportion of false
discoveries \\
Knockoffs & controlled variable selection with FDR \\
Permutation / SEA & small \emph{n}, non-standard models \\
\end{longtable}
}

\begin{Shaded}
\begin{Highlighting}[]
\FunctionTok{p.adjust}\NormalTok{(pvals, }\AttributeTok{method =} \StringTok{"BH"}\NormalTok{)}
\end{Highlighting}
\end{Shaded}

Replication crisis shorthand: pre-register + report effect sizes + share
data. Power, not p-values.

\subsection{TRIPOD-AI + fairness}\label{tripod-ai-fairness}

\begin{itemize}
\tightlist
\item
  TRIPOD-AI extends TRIPOD for prediction models built with ML.
\item
  Report: data provenance, training/evaluation splits, subgroup
  performance, calibration, external validation.
\item
  Fairness: stratify AUC / calibration by sex, ancestry, age bands.
  Disparities in calibration are often larger than in AUC.
\end{itemize}

\subsection{\texorpdfstring{Reproducibility at scale with
\texttt{targets}}{Reproducibility at scale with targets}}\label{reproducibility-at-scale-with-targets}

\begin{Shaded}
\begin{Highlighting}[]
\CommentTok{\# \_targets.R}
\FunctionTok{library}\NormalTok{(targets)}
\FunctionTok{tar\_option\_set}\NormalTok{(}\AttributeTok{packages =} \FunctionTok{c}\NormalTok{(}\StringTok{"tidyverse"}\NormalTok{))}
\FunctionTok{list}\NormalTok{(}
  \FunctionTok{tar\_target}\NormalTok{(raw, }\FunctionTok{read\_csv}\NormalTok{(}\StringTok{"data/raw.csv"}\NormalTok{)),}
  \FunctionTok{tar\_target}\NormalTok{(clean, }\FunctionTok{clean\_data}\NormalTok{(raw)),}
  \FunctionTok{tar\_target}\NormalTok{(fit, }\FunctionTok{fit\_model}\NormalTok{(clean)),}
  \FunctionTok{tar\_target}\NormalTok{(report, quarto}\SpecialCharTok{::}\FunctionTok{quarto\_render}\NormalTok{(}\StringTok{"report.qmd"}\NormalTok{))}
\NormalTok{)}
\end{Highlighting}
\end{Shaded}

\begin{Shaded}
\begin{Highlighting}[]
\ExtensionTok{Rscript} \AttributeTok{{-}e} \StringTok{\textquotesingle{}targets::tar\_make()\textquotesingle{}}
\end{Highlighting}
\end{Shaded}

\subsection{Decision rule}\label{decision-rule-15}

\begin{itemize}
\tightlist
\item
  Bulk DE → DESeq2 for small samples, edgeR/limma for large.
\item
  Many hypotheses → BH by default; knockoffs for variable selection.
\item
  scRNA-seq → Seurat pipeline with batch correction + cluster
  annotation.
\item
  ML prediction model → TRIPOD-AI + subgroup analysis.
\item
  Any multi-step analysis → \texttt{targets} pipeline.
\end{itemize}

\subsection{Common pitfalls}\label{common-pitfalls-15}

\begin{itemize}
\tightlist
\item
  Running DE on a pathway-level aggregate rather than gene-level.
\item
  Cherry-picking a pathway database post-hoc.
\item
  Over-clustering scRNA-seq data until ``known'' groups appear.
\item
  Calling an ML model ``fair'' after checking only overall AUC.
\item
  Submitting a paper without a \texttt{sessionInfo()} or a lock file.
\end{itemize}

\subsection{Further reading}\label{further-reading-15}

\begin{itemize}
\tightlist
\item
  Love, Huber, Anders (2014), \emph{Moderated estimation of fold
  change\ldots{}}
\item
  Collins et al.~(2024), \emph{TRIPOD-AI}.
\item
  Peng, \emph{Reproducible Research with R}.
\end{itemize}

\end{document}
